\documentclass[11pt]{amsart}
\usepackage[utf8]{inputenc}
\usepackage[T1]{fontenc}
\usepackage{lmodern}
\DeclareUnicodeCharacter{2194}{\ensuremath{\leftrightarrow}}
\DeclareUnicodeCharacter{2227}{\ensuremath{\wedge}}
\DeclareUnicodeCharacter{2265}{\ensuremath{\geq}}
\DeclareUnicodeCharacter{2200}{\ensuremath{\forall}}
\DeclareUnicodeCharacter{2208}{\ensuremath{\in}}
\DeclareUnicodeCharacter{2115}{\ensuremath{\mathbb{N}}}
\DeclareUnicodeCharacter{2203}{\ensuremath{\exists}}
\DeclareUnicodeCharacter{2192}{\ensuremath{\to}}
\DeclareUnicodeCharacter{2264}{\ensuremath{\leq}}
\DeclareUnicodeCharacter{1D9C}{\ensuremath{^{\mathrm{c}}}}
\DeclareUnicodeCharacter{2260}{\ensuremath{\neq}}
\DeclareUnicodeCharacter{2209}{\ensuremath{\notin}}
\DeclareUnicodeCharacter{27E8}{\ensuremath{\langle}}
\DeclareUnicodeCharacter{27E9}{\ensuremath{\rangle}}
\DeclareUnicodeCharacter{2228}{\ensuremath{\vee}}
\DeclareUnicodeCharacter{00AC}{\ensuremath{\neg}}
\DeclareUnicodeCharacter{2295}{\ensuremath{\oplus}}
\DeclareUnicodeCharacter{00D7}{\ensuremath{\times}}
\usepackage{amsmath,amssymb,amsthm}
\usepackage[margin=1.1in]{geometry}
\usepackage{enumitem}
\usepackage{booktabs}
\usepackage{longtable}
\usepackage[hidelinks]{hyperref}

\emergencystretch=3em
\newtheorem{theorem}{Theorem}[section]
\newtheorem{lemma}[theorem]{Lemma}
\newtheorem{proposition}[theorem]{Proposition}
\newtheorem{corollary}[theorem]{Corollary}
\theoremstyle{definition}
\newtheorem{definition}[theorem]{Definition}
\theoremstyle{remark}
\newtheorem{remark}[theorem]{Remark}
\newtheorem{example}[theorem]{Example}

\let\oldunderscore\_
\renewcommand{\_}{\oldunderscore\allowbreak}
\newcommand{\lean}[1]{\textup{(Lean: {\small\texttt{#1}})}}
\newcommand{\leanst}[1]{\par{\raggedright\small\textup{(Lean: \texttt{#1})}\par}}
\newcommand{\cC}{\mathcal{C}}
\newcommand{\fl}[1]{\left\lfloor #1\right\rfloor}
\newcommand{\cl}[1]{\left\lceil #1\right\rceil}
\newcommand{\MBU}{\mathrm{MB}}
\newcommand{\RT}{\mathrm{RT}}
\newcommand{\TA}{\mathrm{A}}
\newcommand{\TB}{\mathrm{B}}
\newcommand{\cost}{\operatorname{cost}}
\newcommand{\Ev}{\operatorname{Ev}}
\newcommand{\ind}[1]{[\,#1\,]}

\newcommand{\RE}{\mathrm{RE}}
\newcommand{\pr}{p_r}

\title[A machine-checked reduction of Gallai's conjecture]{A machine-checked reduction of Gallai's
path decomposition conjecture to two rooted tests on $3$-connected graphs}
\author{[Author]}
\date{November 2026}

\begin{document}

\begin{abstract}
Gallai's conjecture asserts that the edges of every connected graph on $n$ vertices can be
partitioned into at most $\lceil n/2\rceil$ paths. We prove that the conjecture is equivalent to two
rooted deletion tests on $3$-connected graphs. Let $J$ be $3$-connected, let $z\in V(J)$, and
suppose that $J-z$ has no matching cut. Test (T1): if $z$ is the only vertex of even degree, then for
every neighbour $r$ of $z$ the graph $J-z$ has a path decomposition into at most $(|V(J)|-1)/2$
paths, at least two of which end at $r$. Test (T2): if $z$ and a neighbour $r$ of $z$ are the only
vertices of even degree, then $J-z$ has a path decomposition into at most $|V(J)|/2$ paths, at least
three of which end at $r$. The proof follows the author's marked-reduction manuscript (Theorems 59,
113 and 114 there) in the simplified form of its formalisation. We include a self-contained proof of
Lov\'asz's theorem on graphs with at most one vertex of even degree, in the strong form in which
every odd vertex is an end of exactly one path, derived from Fan's addible-edge lemma, which we also
prove. On a cone $J=B+z$ over a graph $B$, test (T1) at the apex is exactly the rooted Gallai
statement RE for the Eulerian base: $B$ has a decomposition into at most $|V(B)|/2$ paths with two
ends at any prescribed vertex. Hence the reduction needs RE on $2$-connected Eulerian graphs of
even order without a matching cut, and Gallai's conjecture implies RE. Every result is formalised in
Lean~4 with Mathlib. The equivalence is a single theorem with no hypotheses that depends only on the
standard axioms. The only result we quote without proof is Hall's marriage theorem.
\end{abstract}

\maketitle

\section{Introduction}\label{sec:intro}

All graphs are finite and simple. A \emph{path} is a sequence of distinct vertices $v_0v_1\cdots
v_k$ with $k\ge1$ and $v_{i-1}v_i$ an edge for every $i$; its \emph{ends} are $v_0$ and $v_k$. A
\emph{path decomposition} of a graph $G$ is a family of paths whose edge sets partition $E(G)$, and
the \emph{path number} $p(G)$ is the least size of a path decomposition ($p(G)=0$ when $G$ has no
edges).

\begin{quote}
\textbf{Gallai's conjecture.} If $G$ is connected and $|V(G)|=n$, then $p(G)\le\lceil n/2\rceil$.
\end{quote}

The conjecture is attributed to Gallai and was published by Lov\'asz~\cite{L68}. The bound is
attained: for $k\ge1$, $K_{2k+1}$ has $k(2k+1)$ edges and a path in it has at most $2k$ edges, so
$p(K_{2k+1})\ge k+1=\lceil (2k+1)/2\rceil$. Lov\'asz~\cite{L68} proved the conjecture for graphs
with at most one vertex of even degree (Theorem~\ref{thm:lovasz} below). The parity condition was
weakened by Pyber~\cite{Pyb96} (the even-degree vertices induce a forest) and by
Fan~\cite{Fan05} (each block of the subgraph induced by the even-degree vertices is triangle-free
with maximum degree at most three), and further parity classes were treated by Botler and
Sambinelli~\cite{BS21}. Among other classes, the conjecture is known for graphs of maximum degree at
most five (Bonamy and Perrett~\cite{BP19}) and for graphs of treewidth at most three (Botler,
Sambinelli, Coelho and Lee~\cite{BSCL20}). It is open in general.

\subsection*{The reduction}
We present one chain of reductions, the \emph{spine}, from the author's marked-reduction
manuscript~\cite{MR} (cited as MR). It reduces Gallai's conjecture to two \emph{rooted deletion
tests} on $3$-connected graphs, and it runs in both directions. Recall that $G$ is
\emph{$k$-connected} if $|V(G)|\ge k+1$ and $G-X$ is connected for every $X\subseteq V(G)$ with
$|X|<k$, and that $G$ has a \emph{matching cut} if $V(G)$ can be split into two nonempty sets so
that every vertex has at most one neighbour on the other side.

\begin{definition}[the deletion tests]\label{def:tests}
Let $J$ be a graph, $z\in V(J)$ and $n=|V(J)|$. A vertex is \emph{even} or \emph{odd} according to
the parity of its degree.
\begin{itemize}[leftmargin=3.2em]
\item[(T1)] If $z$ is the only even vertex of $J$, then for every neighbour $r$ of $z$ the graph
$J-z$ has a path decomposition into at most $(n-1)/2$ paths, at least two of which end at $r$.
\item[(T2)] For every neighbour $r$ of $z$: if $z$ and $r$ are the only two even vertices of $J$,
then $J-z$ has a path decomposition into at most $n/2$ paths with at least three path ends at $r$.
\end{itemize}
\end{definition}

A path has two distinct ends, so ``at least three path ends at $r$'' means that at least three of
the paths end at $r$. The number of odd vertices of a graph is even, so in (T1) $n$ is odd and in
(T2) $n$ is even. In both tests the budget is therefore $\lceil |V(J-z)|/2\rceil$, the Gallai bound
for $J-z$.

\begin{theorem}[main theorem; MR Theorem 114]\label{thm:main}
The following are equivalent.
\begin{enumerate}[label=\textup{(\roman*)}]
\item Gallai's conjecture.
\item \textup{(T1)} and \textup{(T2)} hold for every $3$-connected graph $J$ and every
$z\in V(J)$.
\item \textup{(T1)} and \textup{(T2)} hold for every $3$-connected graph $J$ and every
$z\in V(J)$ such that $J-z$ has no matching cut.
\end{enumerate}
\leanst{Gallai.gallai\_iff\_spineTests'}
\end{theorem}

The proof of (iii)$\Rightarrow$(i) uses the tests only at special instances, the \emph{parity
apexes} of Section~\ref{sec:thm114}; Corollary~\ref{cor:mincex} states what a smallest
counterexample looks like.

\subsection*{Why the tests are restrictive}
In a $3$-connected graph every vertex has degree at least three. In (T1) every vertex other than $z$
is odd, and in (T2) every vertex other than $z,r$ is odd. So a test graph is as close as possible to
the graphs covered by Lov\'asz's theorem: $J$ itself has at most two even vertices and, if it has
one, $p(J)\le\fl{n/2}$ by Lov\'asz. Deleting $z$ makes exactly its neighbours change parity. In
(T1) the even vertices of $J-z$ are precisely the neighbours of $z$, and $r$ is one of them; in (T2)
they are the neighbours of $z$ other than $r$, and $r$ becomes odd. Thus the whole difficulty of the
conjecture is concentrated in one deleted vertex and its neighbourhood, inside a graph with no
$2$-separations, no matching cut after the deletion, and odd degrees everywhere else. The
\emph{deletion degree} $d_J(z)$ is at least four (it is even and at least three). The tests are
equivalent to the conjecture, so they are exactly as hard as the conjecture; their value is that
every potential counterexample is forced into this rigid shape.

\subsection*{Lov\'asz's theorem}
Lov\'asz's theorem is the base case of the reduction, and we prove it here in the following strong
form.

\begin{theorem}[Lov\'asz~\cite{L68}, strong form]\label{thm:lovasz-intro}
If at most one vertex of $G$ has positive even degree, then $G$ has a path decomposition in which
every vertex of odd degree is an end of exactly one path and no other vertex is an end. In
particular $p(G)\le\fl{|V(G)|/2}$ whenever $G$ has at most one vertex of even degree.
\leanst{Gallai.lovasz\_ends, Gallai.lovaszStatement}
\end{theorem}

The proof (Section~\ref{sec:lovasz}) is a short induction whose only tool is Fan's addible-edge
lemma~\cite[Lemma~3.4]{Fan05}. We include a complete proof of Fan's lemma by an exchange argument
along chains of paths, so that the paper does not depend on Lov\'asz's path-and-cycle theorem. This
derivation is that of the working note~\cite{LF}. We do not know whether it appears in the
literature.

\subsection*{Cones and rooted Gallai for Eulerian graphs}
A graph is \emph{Eulerian} if all its degrees are even; connectivity is stated separately. For a
graph $B$ on $n$ vertices and $r\in V(B)$ let $\pr(B)$ be the least size of a decomposition of $B$
with at least one path ending at $r$ (when $d_B(r)>0$). If $B$ is Eulerian, every such decomposition has at least two
paths ending at $r$, by parity.

\begin{definition}[rooted Gallai for Eulerian graphs]\label{def:RE}
$\RE(B,r)$: $\pr(B)\le n/2$, for a $2$-connected Eulerian graph $B$ of even order $n$ and $r\in
V(B)$.
\end{definition}

The \emph{cone} $B+z$ adds a vertex $z$ adjacent to every vertex of $B$.

\begin{theorem}[cones; Section~\ref{sec:cones}]\label{thm:cones-intro}
Let $B$ be a graph of order $n$ and $J=B+z$.
\begin{enumerate}[label=\textup{(\alph*)}]
\item $J$ is $3$-connected if and only if $B$ is $2$-connected, and $J-z$ has a matching cut if and
only if $B$ has one.
\item \textup{(T1)} holds at $(J,z)$ if and only if: when $n$ is even and $B$ is Eulerian, $B$
has, for every $r\in V(B)$, a decomposition into at most $n/2$ paths with at least two ends at $r$.
\textup{(T2)} always holds at $(J,z)$.
\item Assertion \textup{(iii)} of Theorem~\ref{thm:main} implies $\RE(B,r)$ for every $2$-connected
Eulerian graph $B$ of even order without a matching cut and every $r$. Conversely, for $B$ of even
order, \textup{(T1)} and \textup{(T2)} hold at every vertex of $B+z$ if and only if $\RE(B,r)$ holds
for every $r$ whenever $B$ is Eulerian.
\item Gallai's conjecture implies $\pr(B)\le n/2$ for every connected Eulerian graph $B$ of even
order $n$ and every vertex $r$ of positive degree; in particular it implies RE.
\end{enumerate}
\leanst{Gallai.kConnected\_three\_fullCone\_iff, Gallai.hasMatchingCut\_delVert\_fullCone\_iff,
Gallai.test1\_fullCone\_iff, Gallai.test2\_fullCone\_apex,
Gallai.rootedDecomp\_of\_spineTestsNMC, Gallai.fullConeTests\_iff,
Gallai.rootedDecomp\_of\_gallai}
\end{theorem}

So, at the apex of a cone, the spine asks exactly for RE on the $2$-connected Eulerian graphs of even
order without a matching cut, and
\[
\text{Gallai}\ \Longrightarrow\ \RE\ \Longrightarrow\ \text{Gallai for $2$-connected Eulerian
graphs of even order.}
\]
This correspondence is taken from the working note~\cite{EC}, Section~2, where cones were first
singled out as a test case.

\subsection*{The chain of reductions}
The proof of Theorem~\ref{thm:main} has six steps.
\begin{enumerate}[label=\arabic*.,leftmargin=2em]
\item \emph{Lov\'asz's theorem} (Section~\ref{sec:lovasz}), from Fan's lemma.
\item \emph{Marked path numbers} (Section~\ref{sec:marked}; MR Theorem~1). Each mark in a set $S$ may
absorb one path ending there; $q(F,S)$ counts the paths that are not absorbed. The marked bound
$q(F,S)\le\fl{(|V(F)|-|S|)/2}$ implies Gallai's bound with one mark. On graphs whose $2$-connected
blocks are small, it follows from the marked bound on $2$-connected graphs and a three-end condition.
\item \emph{Block transfer and matching cuts} (Section~\ref{sec:transfer}; MR \S2, Lemma~6,
Theorem~7). The block transfer is proved by induction through leaves and cut vertices. A
$2$-connected graph with a matching cut satisfies the marked bound once smaller blocks do.
\item \emph{Theorem 59} (Section~\ref{sec:thm59}). Two rooted assertions on $2$-connected graphs
without a matching cut imply the marked bound everywhere, order by order.
\item \emph{Separating pairs} (Section~\ref{sec:sep}; MR Theorem~113). A component behind a
separating pair whose vertices are all odd can be cut off, which reduces the rooted assertions to
smaller graphs.
\item \emph{Deletion tests} (Section~\ref{sec:thm114}; MR Theorem~114). At a smallest failure of the
rooted assertions, adding a vertex $z$ adjacent to all even vertices and the root gives a
$3$-connected graph on which a deletion test fails.
\end{enumerate}
Section~\ref{sec:cones} treats cones, Section~\ref{sec:use} explains how the reduction is meant to
be used, and Section~\ref{sec:lean} describes the formalisation.

\subsection*{Status of the results}
Every numbered theorem, proposition, lemma and corollary is proved in the text. The only result
quoted without proof is Hall's marriage theorem~\cite{Hall35}, used once
(Lemma~\ref{lem:hall}). There are no computer-assisted claims. Every result is also formalised in
Lean, with the few exceptions listed in Section~\ref{sec:lean}, which are formalised only in the
special form the main theorem uses or not at all. The formalisation is an independent check and is
not needed to read the paper. Conjectures are labelled as such.

\subsection*{Formalisation and relation to MR}
The whole chain is formalised in Lean~4~\cite{Lean4} with Mathlib~\cite{Mathlib}. The theorem
\texttt{Gallai.gallai\_iff\_spineTests'} states Theorem~\ref{thm:main}. It has no hypotheses and
depends only on the axioms \texttt{propext}, \texttt{Classical.choice} and \texttt{Quot.sound}
(Section~\ref{sec:lean}). The formal proof follows MR, but in several places it is organised
differently or more simply, and we follow the formal proof at each of them. The project records
these departures in its audit reports (A14 and A15), in the module documentation of
\texttt{MarkedBlocks.lean}, in the proofs in \texttt{ThreeConnected.lean}, and in the section
``Deviations from the manuscript proof'' of \texttt{lean/README.md}:
\begin{itemize}
\item a single assignment lemma proved with Hall's theorem (Lemma~\ref{lem:hall}) replaces two
arguments of MR~\S2.1, the private-leaf identity and the odd-mark rank inequality (A14);
\item the bridge formula of MR~\S2 is not needed: a bridge whose sides both have at least two
vertices gives a cut vertex with two large sides, or else ends in a leaf (\texttt{MarkedBlocks.lean});
\item the matching cut with connected shores is obtained from components
(Lemma~\ref{lem:shores}) rather than from an inclusion-minimal disconnecting submatching (A14);
\item the induction of Theorem~59 runs explicitly on the bound for the order of the
$2$-connected blocks (Theorem~\ref{thm:59}); MR's phrase ``induct on the order of a two-connected
block'' says this, and the formal statement makes it explicit (A14);
\item singletons of odd order are handled by the two-mark cone and forgetting one mark
(Proposition~\ref{prop:classC}) instead of adjoining a degree-two vertex and trimming (A14);
\item the matching-cut step needs only connectivity (Theorem~\ref{thm:mc}; A14);
\item the forward direction of Theorem~\ref{thm:main} is proved directly from the bridge double
(Lemma~\ref{lem:gallai-rooted}); the bridge and diamond identities of MR~\S17 are not used
(see the proofs of \texttt{test1\_of\_gallai} and \texttt{test2\_of\_gallai});
\item in the separator theorem a single root analysis (Lemma~\ref{lem:root}) serves all cases, and
the four sub-cases of the even/even case become at most two grafts (Lemma~\ref{lem:lift-ee};
\texttt{lean/README.md}).
\end{itemize}
The project's audits A1 and A2 of MR found no error but recorded gaps in its exposition. We use the
corrected statements throughout. Theorem~59 is stated with an explicit marked-bound conclusion (A2,~F1).
Theorem~113 is stated, as A2 (F7) recommends, without minimality and without excluding matching
cuts; its second hypothesis is the marked bound $(\MBU_{n-2})$, which F7's hypothesis (both tests
below order $n$) implies by Theorem~59. The marked tests on blocks
without a matching cut are derived for every nonempty marking, not only for even marks (A1, G1 and
G7). The instances of order $m+1$ used by Theorem~59 are named (A1,~G10). Only the two-part case of
MR Theorem~7, a matching cut, is used (A1,~G6). Single-vertex deletions are covered in the apex lemma
(A2,~F6). MR Theorem~57 and the doubling identities of MR~\S17 (A1, G8 and G9) are not needed.

The present paper supersedes three working notes of the project for citation: the spine
paper~\cite{SP}, which contains Sections~\ref{sec:prelim} and~\ref{sec:marked}--\ref{sec:thm114}
in an earlier form; the note~\cite{LF} on Lov\'asz's theorem; and Section~2 of~\cite{EC}. None of
the arguments here depends on those notes or on MR: every proof is given in full.

\section{Preliminaries}\label{sec:prelim}

\subsection{Graphs}
For a graph $G$ we write $V(G)$, $E(G)$, $d_G(v)$ and $N_G(v)$, and call $|V(G)|$ the
\emph{order} of $G$. A vertex $v$ is \emph{even} or \emph{odd} according to the parity of $d_G(v)$;
a vertex of degree $0$ is even. $\Ev(G)$ is the set of even vertices. By the handshake lemma the
number of odd vertices is even. For $U\subseteq V(G)$, $G[U]$ is the induced subgraph and
$G-X=G[V(G)\setminus X]$. A \emph{leaf} is a vertex of degree one.

$G$ is \emph{$k$-connected} if $|V(G)|\ge k+1$ and $G-X$ is connected for every $X\subseteq V(G)$
with $|X|<k$ \lean{Gallai.KConnected}. Thus $G$ is $2$-connected if and only if it has at least three
vertices, is connected, and has no \emph{cut vertex}, that is, no $x$ with $G-x$ disconnected. (The
block transfer in Lean states its tests with this no-cut-vertex form, \texttt{TwoConnected}, and
converts it with \texttt{Gallai.TwoConnected.kConnected}.) If
$x$ is a cut vertex of a connected graph $G$, then $V(G)\setminus\{x\}=C\sqcup D$ with $C,D$
nonempty and no edge between $C$ and $D$; we call $(C,D)$ a \emph{separation at $x$}. We say that
the $2$-connected blocks of $G$ have order at most $m$ if $|U|\le m$ for every $U\subseteq V(G)$ such
that $G[U]$ is $2$-connected \lean{Gallai.SmallBlocks}.

\begin{lemma}\label{lem:kconn}
Let $G$ be $k$-connected. Then $d_G(v)\ge k$ for every $v$; $G$ is connected if $k\ge1$; and
$G-z$ is connected for every vertex $z$ if $k\ge2$.
\leanst{Gallai.KConnected.le\_degree, Gallai.KConnected.connected,
Gallai.KConnected.delVert\_connected}
\end{lemma}

\begin{proof}
If $d_G(v)<k$, then $|N_G(v)\cup\{v\}|\le k<|V(G)|$, so some $w$ is neither $v$ nor a neighbour of
$v$; but $v$ is isolated in $G-N_G(v)$, which has fewer than $k$ deleted vertices and contains $w$, a
contradiction. The other two claims are the cases $X=\emptyset$ and $X=\{z\}$ of the definition.
\end{proof}

Three constructions recur. The graph $G^x$ is $G$ with a new vertex (a \emph{private leaf}) adjacent
only to $x$ \lean{Gallai.addLeaf}. For $S\subseteq V(G)$, the \emph{cone} $G+z(S)$ is $G$ with a new
vertex $z$ adjacent exactly to the vertices of $S$ \lean{Gallai.apex, Gallai.cone}. For
$z\in V(J)$, $J-z$ is the subgraph induced on $V(J)\setminus\{z\}$
\lean{Gallai.delVert}.

A \emph{matching cut} of $G$ is a partition $V(G)=X\sqcup Y$ with $X,Y\ne\emptyset$ such that every
vertex has at most one neighbour on the other side; equivalently, the edges between $X$ and $Y$
form a (possibly empty) matching \lean{Gallai.HasMatchingCut}. The class $\cC$ consists of the
$2$-connected graphs without a matching cut \lean{Gallai.InClassC}.

\subsection{Decompositions and ends}
Paths are as in Section~\ref{sec:intro}; a path and its reversal are the same path. We write $V(P)$
and $E(P)$ for the vertices and edges of a path $P$, and regard a decomposition $D$ as a list of
paths, so $|D|$ is its number of paths \lean{Gallai.PathDecomp}. Every graph has a decomposition
(its edges, one per path), so $p(G)$ is attained. For a path $P$ and a vertex $v$ let
$b_P(v)\in\{0,1\}$ be the number of ends of $P$ at $v$. For a family $D$ of paths let
$b_D(v)=\sum_{P\in D}b_P(v)$, the number of paths of $D$ ending at $v$ \lean{Gallai.PathDecomp.ends}.
Put $\ind{\phi}=1$ if $\phi$ holds and $0$ otherwise.

\begin{lemma}[parity]\label{lem:parity}
For every decomposition $D$ of $G$ and every vertex $v$:
$\sum_{u}b_D(u)=2|D|$, $b_D(v)\le d_G(v)$, and $b_D(v)\equiv d_G(v)\pmod 2$.
\leanst{Gallai.PathDecomp.sum\_ends, Gallai.PathDecomp.ends\_mod\_two,
Gallai.PathDecomp.ends\_le\_degree}
\end{lemma}

\begin{proof}
Each path has exactly two ends. A path $P$ through $v$ contains two edges at $v$ if $v$ is internal
and one if $v$ is an end, so $P$ contains $2\ind{v\in V(P)}-b_P(v)$ edges at $v$. The edges at $v$
are partitioned among the paths, so $d_G(v)=2\,|\{P\in D:v\in V(P)\}|-b_D(v)$. This gives the
parity, and $b_D(v)\le |\{P:v\in V(P)\}|\le d_G(v)$ because every path through $v$ contains an edge
at $v$.
\end{proof}

\begin{lemma}[splitting and joining]\label{lem:split}
Let $D$ be a decomposition of $G$.
\begin{enumerate}[label=\textup{(\alph*)}]
\item If $v$ is an internal vertex of $P\in D$, replacing $P$ by its two subpaths ending at $v$ gives
a decomposition $D'$ with $|D'|=|D|+1$, $b_{D'}(v)=b_D(v)+2$ and $b_{D'}(u)=b_D(u)$ for $u\ne v$.
\item If $b_D(v)<d_G(v)$, then $v$ is an internal vertex of some path of $D$.
\item If $P,Q$ are paths ending at $v$ with $V(P)\cap V(Q)=\{v\}$, then $P\cup Q$ is a path whose ends
are the other ends of $P$ and $Q$.
\end{enumerate}
\leanst{Gallai.PathDecomp.exists\_split, Gallai.PathDecomp.exists\_internal\_of\_ends\_lt,
Gallai.GPath.exists\_join}
\end{lemma}

\begin{proof}
(a) and (c) are immediate. For (b), if $v$ is internal to no path, every path through $v$ ends at
$v$ and contains exactly one edge at $v$, so $d_G(v)=b_D(v)$.
\end{proof}

\begin{lemma}[two and three ends]\label{lem:ends}
Let $D$ be a decomposition of $G$, $x\in V(G)$ and $k\ge0$.
\begin{enumerate}[label=\textup{(\alph*)}]
\item If $d_G(x)$ is even and positive and $|D|+\ind{b_D(x)=0}\le k$, then some decomposition $D'$
has $|D'|\le k$ and $b_{D'}(x)\ge2$.
\item If $d_G(x)$ is odd and at least $3$ and $|D|\le k$, then some decomposition $D'$ has
$|D'|\le k+1$ and $b_{D'}(x)\ge3$.
\end{enumerate}
\leanst{Gallai.exists\_two\_ends, Gallai.exists\_three\_ends}
\end{lemma}

\begin{proof}
(a) If $b_D(x)=0<d_G(x)$, split at an internal visit of $x$ (Lemma~\ref{lem:split}); the result has
$|D|+1\le k$ paths and two ends at $x$. Otherwise $b_D(x)\ge1$ is even by Lemma~\ref{lem:parity},
so $b_D(x)\ge2$ and $|D|\le k$. (b) If $b_D(x)\ge3$ take $D'=D$. Otherwise $b_D(x)=1$ by parity, and
$1<3\le d_G(x)$, so we may split at an internal visit of $x$.
\end{proof}

For a decomposition $D$ of $G$ and a vertex $x$ put $\lambda_D(x)=|D|+\ind{b_D(x)=0}$
\lean{Gallai.leafCost}.

\begin{lemma}[private leaves]\label{lem:leaf}
Let $x\in V(G)$ and let $\ell$ be the private leaf of $G^x$.
\begin{enumerate}[label=\textup{(\alph*)}]
\item Every decomposition $D'$ of $G^x$ yields a decomposition $D$ of $G$ with either
$|D|\le|D'|$ and $b_D(x)=b_{D'}(x)+1$, or $|D|\le|D'|-1$ and $b_D(x)=b_{D'}(x)-1$.
\item For every decomposition $D$ of $G$, $p(G^x)\le\lambda_D(x)$.
\item Some decomposition $D$ of $G$ has $\lambda_D(x)\le p(G^x)$.
\end{enumerate}
\leanst{Gallai.PathDecomp.trim\_ends, Gallai.pathNumber\_addLeaf\_le, Gallai.exists\_leafCost\_le}
\end{lemma}

\begin{proof}
(a) The edge $x\ell$ lies on one path $Q$, and $\ell$ is an end of $Q$. If $Q=x\ell$, delete $Q$:
one path fewer, and $x$ loses the end of $Q$. Otherwise delete $\ell$ from $Q$: the new path ends at
$x$, where $Q$ did not end. (b) If $b_D(x)\ge1$, extend one path ending at $x$ by the edge $x\ell$;
otherwise add $x\ell$ as a new path. (c) Apply (a) to a decomposition $D'$ of size $p(G^x)$. In the
first case $b_D(x)\ge1$ and $\lambda_D(x)=|D|\le p(G^x)$; in the second
$\lambda_D(x)\le|D|+1\le p(G^x)$.
\end{proof}

\section{Lov\'asz's theorem from Fan's addible-edge lemma}\label{sec:lovasz}

Lov\'asz~\cite{L68} proved his theorem through a more general statement on decompositions into
paths and cycles. We derive it, in the strong form of Theorem~\ref{thm:lovasz-intro}, from Fan's
addible-edge lemma~\cite[Lemma~3.4]{Fan05}. To keep the paper self-contained we prove Fan's lemma
too (Lemmas~\ref{lem:chain-exec}--\ref{lem:fan}). The proof is an exchange argument along chains of
paths; it does not use Lov\'asz's theorem, so there is no circularity. It is the proof formalised in
\texttt{Fan.lean} and \texttt{FanMoves.lean}.

\subsection{Fan's lemma}
Throughout this subsection $K$ is a graph, $D$ a decomposition of $K$, and $u\in V(K)$. A path
of $D$ has two distinct ends, so $b_D(x)$ is the number of paths of $D$ with an end at $x$.

\begin{definition}[predecessors and chains]\label{def:chain}
Let $x\ne u$ and let $P\in D$ be a path with an end at $x$ and with $u\in V(P)$. Traverse $P$
starting from $x$; the vertex visited immediately before $u$ is the \emph{predecessor of $x$ via
$P$}. It is a neighbour of $u$, since consecutive vertices of $P$ are adjacent. A \emph{chain from
$t$} is a sequence $x_0=t,x_1,\dots,x_k$ ($k\ge0$) of distinct vertices, all different from $u$, such
that $x_i$ is the predecessor of $x_{i-1}$ via some path of $D$ for $1\le i\le k$, and some path of
$D$ with an end at $x_k$ does not contain $u$.
\leanst{Gallai.FanPred, Gallai.FanAvoid, Gallai.FanChain}
\end{definition}

\begin{lemma}[executing a chain]\label{lem:chain-exec}
Let $t\ne u$ with $tu\notin E(K)$, and suppose that a chain from $t$ exists. Then $K+tu$ has a
decomposition $D^+$ with $|D^+|=|D|$, $b_{D^+}(t)=b_D(t)-1$, $b_{D^+}(u)=b_D(u)+1$ and
$b_{D^+}(w)=b_D(w)$ for every other vertex $w$.
\leanst{Gallai.fan\_exec, Gallai.PathDecomp.append\_move, Gallai.PathDecomp.reroute\_move,
Gallai.fanChain\_transport}
\end{lemma}

\begin{proof}
Induction on the length $k$ of the chain, for all $K$, $D$ and $t$ at once.

If $k=0$, some path $P\in D$ has an end at $t$ and avoids $u$. Extending $P$ at $t$ by the edge $tu$
gives a path, since $u\notin V(P)$. Replacing $P$ by it gives $D^+$: the end of $P$ at $t$ moves to
$u$, and nothing else changes.

Let $k\ge1$. Some path $P\in D$ has an end at $t$, contains $u$, and reads
$P=t\cdots x_1\,u\cdots s$ from $t$, where $s\ne t$ is its other end. Let $P'$ be the walk that
traverses the segment of $P$ from $x_1$ back to $t$, then the edge $tu$, and then the rest of $P$
from $u$ to $s$. It has the same vertices as $P$, each once, so it is a path, with ends $x_1$ and
$s$. Its edges are those of $P$ with $x_1u$ replaced by $tu$. Hence $D_1=D-P+P'$ is a decomposition
of $K_1=K-x_1u+tu$ with $|D_1|=|D|$, $b_{D_1}(t)=b_D(t)-1$, $b_{D_1}(x_1)=b_D(x_1)+1$, and all
other end counts unchanged. In $K_1$ the vertex $x_1$ is not adjacent to $u$, and $x_1\ne u$.

We claim that $x_1,\dots,x_k$ is a chain from $x_1$ in $D_1$. All paths of $D$ other than $P$ are
paths of $D_1$. A path used in the chain at a vertex $x_i$ with $i\ge1$ has an end at $x_i\ne t$. If
that path is $P$, then $x_i=s$, and $P$ is used to read a predecessor of $s$. Read from $s$, the
path $P$ is $s\cdots w\,u\,x_1\cdots t$, and $P'$ is $s\cdots w\,u\,t\cdots x_1$. Both agree up to
$u$, so the predecessor of $s$ is $w$ in both. Finally, the path with an end at $x_k$ avoiding $u$ is
not $P$, since $P$ contains $u$. This proves the claim.

By induction, $K_1+x_1u=K+tu$ has a decomposition $D^+$ of size $|D_1|=|D|$ in which an end has moved
from $x_1$ to $u$ relative to $D_1$. Combined with the first step, an end has moved from $t$ to $u$.
\end{proof}

\begin{lemma}[existence of a chain]\label{lem:chain-exists}
Let $t\ne u$ with $tu\notin E(K)$, and put $Z=\{w\in N_K(u): b_D(w)=0\}$. If $b_D(t)>|Z|$, then a
chain from $t$ exists.
\leanst{Gallai.fan\_exists\_chain}
\end{lemma}

\begin{proof}
Suppose not. Let $R$ be the set of vertices $x$ for which there is a sequence $t=x_0,\dots,x_j=x$
of distinct vertices with $x_i$ the predecessor of $x_{i-1}$ via a path of $D$ for $1\le i\le j$.
Then $t\in R$, and:
\begin{enumerate}[label=\textup{(\arabic*)}]
\item $R$ is closed under predecessors: if $x\in R$ is reached by the sequence $L$ and $z$ is a
predecessor of $x$, then either $z$ lies on $L$ and is reached by a prefix of $L$, or $L,z$ reaches
$z$.
\item Every path of $D$ with an end in $R$ contains $u$, since otherwise a chain would exist.
\item Every vertex of $R\setminus\{t\}$ is a predecessor, hence a neighbour of $u$. So $u\notin R$
(a vertex is not its own neighbour, and $t\ne u$), and $t\notin N_K(u)$.
\end{enumerate}
For a path $P\in D$ let $e_P$ be the number of ends of $P$ in $R$ and $f_P$ the number of edges of
$P$ joining $u$ to a vertex of $R$. We claim $e_P\le f_P$. If $P$ has an end $x\in R$, then $u\in
V(P)$ by (2) and $u\ne x$ by (3); the predecessor $z$ of $x$ via $P$ lies in $R$ by (1), and $zu$ is
an edge of $P$ on the $x$-side of $u$. If both ends of $P$ lie in $R$, then $u$ is an internal
vertex of $P$, and the two edges obtained, one on each side of $u$, are distinct. This proves the
claim. Each edge of $K$ lies on exactly one path of $D$, so
\[
\sum_{x\in R}b_D(x)=\sum_{P\in D}e_P\le\sum_{P\in D}f_P=|N_K(u)\cap R|.
\]
Every $w\in N_K(u)\cap R$ either lies in $Z$ or has $b_D(w)\ge1$, and $t\notin N_K(u)$. Hence
$|N_K(u)\cap R|\le|Z|+\sum_{x\in R\setminus\{t\}}b_D(x)$. Comparing with the display gives
$b_D(t)\le|Z|$, a contradiction.
\end{proof}

\begin{lemma}[Fan~{\cite[Lemma~3.4]{Fan05}}]\label{lem:fan}
Let $uv\in E(G)$, let $G'=G-uv$, and let $D'$ be a decomposition of $G'$. If
\[
b_{D'}(v)>\bigl|\{w\in N_{G'}(u): b_{D'}(w)=0\}\bigr|,
\]
then $G$ has a decomposition $D$ with $|D|=|D'|$, $b_D(u)=b_{D'}(u)+1$, $b_D(v)=b_{D'}(v)-1$ and
$b_D(w)=b_{D'}(w)$ for every other vertex $w$.
\leanst{Gallai.fan\_single}
\end{lemma}

\begin{proof}
Apply Lemma~\ref{lem:chain-exists} with $K=G'$ and $t=v$; here $v\ne u$ and $uv\notin E(G')$. Then
apply Lemma~\ref{lem:chain-exec}; note $G'+uv=G$.
\end{proof}

The statement agrees with Fan's Lemma~3.4 as it is quoted in the literature, for instance
in~\cite{BS21}.

\subsection{Lov\'asz's theorem}

\begin{theorem}[Lov\'asz, strong form; Theorem~\ref{thm:lovasz-intro}]\label{thm:lovasz-strong}
If at most one vertex of $G$ has positive even degree, then $G$ has a decomposition $D$ with
$b_D(v)=d_G(v)\bmod 2$ for every vertex $v$.
\leanst{Gallai.lovasz\_ends}
\end{theorem}

\begin{proof}
Strong induction on $\nu(G)=\sum_v d_G(v)+\ind{G\text{ has no vertex of positive even degree}}$,
over all finite graphs. Deleting an edge $xy$ lowers the degrees of $x$ and $y$ by one and changes
no other degree. In each case below the claimed decomposition has the right end counts at every
vertex; we check the vertices whose counts change.

\emph{Case I: some vertex $z$ has positive even degree.} Choose a neighbour $u$ of $z$ and put
$G'=G-uz$. In $G'$ the vertex $z$ is odd. The vertex $u$ has positive degree and $u\ne z$, so it is
odd in $G$, and in $G'$ it is even (possibly isolated). Every other vertex keeps its parity. So $u$
is the only possible vertex of positive even degree in $G'$, and $\nu(G')\le\nu(G)-1$. By induction
$G'$ has a decomposition $D'$ with $b_{D'}=d_{G'}\bmod2$; thus $b_{D'}(u)=0$ and $b_{D'}(z)=1$.
Every neighbour of $u$ in $G'$ has positive degree and differs from $u$, so it is odd and has an
end. The right-hand side in Lemma~\ref{lem:fan}, applied to the edge $uz$ with $v=z$, is therefore
$0<1=b_{D'}(z)$. The lemma adds $uz$ and moves an end from $z$ to $u$. The result has $b(u)=1$ and
$b(z)=0$, matching the parities in $G$.

\emph{Case II: every vertex of positive degree is odd, and some $x$ has $d_G(x)\ge3$.} Choose a
neighbour $u$ of $x$, and let $H$ be $G-ux$ with a new leaf $\ell$ attached at $u$. In $H$ the vertex
$u$ keeps its degree, $\ell$ has degree $1$, and $x$ has positive even degree $d_G(x)-1\ge2$. So $x$
is the unique vertex of positive even degree in $H$, and $\nu(H)=\nu(G)-1$: the degree sum is the
same and the indicator drops. By induction $H$ has a decomposition $D_H$ with $b_{D_H}=d_H\bmod2$;
in particular $b(x)=0$, $b(u)=1$ and $b(\ell)=1$. The leaf is an end of its path. Deleting $\ell$
gives a decomposition $D'$ of $G'=G-ux$: if the leaf path was the single edge $u\ell$, it disappears
together with its end at $u$; otherwise the shortened path gains an end at $u$. So $b_{D'}(u)\in
\{0,2\}$, and all other end counts are unchanged.
\begin{itemize}
\item If $b_{D'}(u)=0$, add $ux$ as a new one-edge path.
\item If $b_{D'}(u)=2$, every neighbour of $x$ in $G'$ has an end: $u$ has two, and the others are
odd. Lemma~\ref{lem:fan}, applied to the edge $xu$ with the roles $(u,v)$ played by $(x,u)$, adds
$ux$ and moves an end from $u$ to $x$.
\end{itemize}
In both cases the result has $b(u)=b(x)=1$.

\emph{Case III: every vertex has degree at most $1$, and there is an edge $ux$.} Then $G-ux$ again
has no vertex of positive even degree and $\nu(G-ux)=\nu(G)-2$. Decompose it by induction and add
$ux$ as a one-edge path.

\emph{Case IV: $G$ has no edges.} The empty decomposition works.

The cases are exhaustive. If no vertex has positive even degree and no vertex has degree at least
$3$, every degree is $0$ or $1$.
\end{proof}

\begin{theorem}[Lov\'asz~\cite{L68}]\label{thm:lovasz}
If a graph $G$ on $N$ vertices has at most one vertex of even degree (vertices of degree $0$ count as
even), then $p(G)\le\fl{N/2}$.
\leanst{Gallai.lovaszStatement}
\end{theorem}

\begin{proof}
At most one vertex has positive even degree, so Theorem~\ref{thm:lovasz-strong} gives $D$ with
$2|D|=\sum_v b_D(v)=o$, the number of odd vertices (Lemma~\ref{lem:parity}). Since $o$ is even and
$o\le N$, we get $|D|=o/2\le\fl{N/2}$.
\end{proof}

Besides the statement itself we use the following tight form.

\begin{corollary}[tight form]\label{cor:lovasz}
Let $G$ be a graph and $U\subseteq V(G)$ a set containing both ends of every edge. If at most one
vertex of $U$ is even, then $G$ has a decomposition $D$ with $2|D|=|\{v:d_G(v)\text{ odd}\}|$ and
$b_D(v)=\ind{d_G(v)\text{ odd}}$ for every $v$. In particular $|D|=\fl{|U|/2}$.
\leanst{Gallai.lovasz\_tight, Gallai.lovasz\_on}
\end{corollary}

\begin{proof}
$G[U]$ has the same edges and degrees on $U$, so Theorem~\ref{thm:lovasz} gives a decomposition $D$
with $2|D|\le 2\fl{|U|/2}$. All odd vertices lie in $U$; let $o$ be their number, which is even. If
$U$ has no even vertex then $o=|U|\ge 2\fl{|U|/2}$; if it has one then $|U|=o+1$ is odd and
$2\fl{|U|/2}=o$. In both cases $2|D|\le o$. On the other hand $2|D|=\sum_v b_D(v)\ge o$, since every
odd vertex has $b_D\ge1$ by parity. Hence equality holds throughout: every odd vertex is an end of
exactly one path and no other vertex is an end.
\end{proof}

\section{Marked path numbers and the marked tests}\label{sec:marked}

This section follows MR~\S1. Its main result, Theorem~\ref{thm:transfer}, is the sufficiency half
of MR Theorem~1 in the finite-order form needed for Theorem~59; it is proved in
Section~\ref{sec:transfer}.

\subsection{The marked number}

\begin{definition}[marked partitions]\label{def:marked}
Let $F$ be a graph and $S\subseteq V(F)$ a set of \emph{marks}. A \emph{marked partition} of
$(F,S)$ is a pair $(D,\sigma)$ where $D$ is a decomposition of $F$ and $\sigma$ is an injective
partial map from the paths of $D$ to $S$ such that $\sigma(P)$ is an end of $P$ whenever $\sigma(P)$
is defined. A path in the domain of $\sigma$ is \emph{assigned} (to the mark $\sigma(P)$). The
\emph{cost} of $(D,\sigma)$ is the number of unassigned paths, and the \emph{marked number}
$q(F,S)$ is the least cost of a marked partition of $(F,S)$. The \emph{marked bound} for $(F,S)$ is
\begin{equation}\label{eq:mb}
q(F,S)\le\fl{(|V(F)|-|S|)/2}.
\end{equation}
\leanst{Gallai.MarkedPartition, Gallai.markedNumber}
\end{definition}

Thus $q(F,\emptyset)=p(F)$. Since $\sigma$ is injective, a mark absorbs at most one path, and since
$\sigma$ is a map, each path is absorbed by at most one mark.

\begin{lemma}\label{lem:marked-basic}
Let $F$ be a graph, $S\subseteq V(F)$ and $x,y\in V(F)$.
\begin{enumerate}[label=\textup{(\alph*)}]
\item $p(F)\le q(F,S)+|S|$.
\item If $S'\subseteq S$ then $q(F,S)\le q(F,S')$.
\item $q(F,S)\le q(F,S\cup\{y\})+1$.
\item If some decomposition $D$ has $b_D(x)\ge1$, then $q(F,\{x\})\le|D|-1$.
\item Some decomposition $D$ of $F$ has $\lambda_D(x)\le q(F,\{x\})+1$.
\end{enumerate}
\leanst{Gallai.pathNumber\_le\_markedNumber\_add, Gallai.markedNumber\_anti,
Gallai.markedNumber\_le\_insert, Gallai.markedNumber\_singleton\_add\_one,
Gallai.exists\_leafCost\_le\_markedNumber}
\end{lemma}

\begin{proof}
(a) In a marked partition of cost $c$ at most $|S|$ paths are assigned, so $|D|\le c+|S|$.
(b) A marked partition of $(F,S')$ is one of $(F,S)$. (c) In an optimal marked partition of
$(F,S\cup\{y\})$, unassign the path assigned to $y$, if any. (d) Assign one path ending at $x$ to
$x$. (e) Let $(D,\sigma)$ be optimal for $(F,\{x\})$. If some path is assigned to $x$, then
$b_D(x)\ge1$ and $\lambda_D(x)=|D|=q(F,\{x\})+1$. Otherwise $\lambda_D(x)\le|D|+1=q(F,\{x\})+1$.
\end{proof}

By (a), the marked bound with a single mark gives $p(F)\le\fl{(|V(F)|-1)/2}+1=\cl{|V(F)|/2}$, which
is Gallai's bound. MR shows conversely that Gallai's conjecture implies the marked bound for every
connected $F$ and nonempty $S$ (the doubled cone $W(F,S)$ of MR~\S1; Lean:
\texttt{Gallai.gallai\_iff\_markedBound}). We do not need this converse.

\subsection{The cone inequality}

\begin{lemma}[cone inequality; MR \S1]\label{lem:cone}
Let $X=F+z(S)$. Every decomposition $D$ of $X$ satisfies
\[
|S|+2q(F,S)+b_D(z)\le 2|D|.
\]
\leanst{Gallai.le\_outEnds, Gallai.apex\_outEnds\_bound}
\end{lemma}

\begin{proof}
A path contains $z$ at most once. Split every path of $D$ at $z$ whenever $z$ is one of its internal
vertices (Lemma~\ref{lem:split}(a)). The result $D^*$ has $b_{D^*}(v)=b_D(v)$ for every $v\ne z$, and
in $D^*$ every path through $z$ ends at $z$ and so contains exactly one edge $zs$ ($s\in S$). Every
edge $zs$ lies on exactly one path of $D^*$, and that path contains no other edge at $z$. Hence the
paths of $D^*$ through $z$ correspond bijectively to $S$ via their edge at $z$. For such a path $P$
with edge $zs$: if $P\ne zs$, then $P-z$ is a path of $F$ with end $s$, and we assign it to $s$; if
$P=zs$ we discard it. The paths of $D^*$ avoiding $z$ are paths of $F$ and stay unassigned. This is
a marked partition of $(F,S)$: its edges are those of $X$ except the edges at $z$, and distinct paths
receive distinct marks. Its cost is the number $a$ of paths of $D^*$ avoiding $z$, so $q(F,S)\le a$.
Every path of $D^*$ through $z$ has exactly one end other than $z$, and every other path has two, so
\[
2|D|-b_D(z)=\sum_{v\ne z}b_D(v)=\sum_{v\ne z}b_{D^*}(v)=|S|+2a\ge|S|+2q(F,S).\qedhere
\]
\end{proof}

(In fact $\min_D(2|D|-b_D(z))=|S|+2q(F,S)$, the identity of MR~\S1 \lean{Gallai.minOutEnds\_eq}.
Only the inequality is used.)

\subsection{The assignment lemma}
The following lemma replaces both the private-leaf identity and the odd-mark rank inequality of
MR~\S2.1.

\begin{lemma}[assignment lemma]\label{lem:hall}
Let $D$ be a decomposition of $F$ and let $T,O\subseteq V(F)$. Suppose $b_D(v)\ge1$ for every
$v\in T$ and $b_D(v)\ge2$ for every $v\in T\setminus O$. Then
\[
q(F,T)+|T|\le|D|+\fl{|O|/2}.
\]
\leanst{Gallai.markedNumber\_add\_card\_le}
\end{lemma}

\begin{proof}
Let $k=\fl{|O|/2}$. Form a bipartite graph between $T$ and the set of slots consisting of the paths
of $D$ together with $k$ dummy slots: a mark $v$ is adjacent to every path of $D$ ending at $v$ and to
every dummy slot. We check Hall's condition. Let $s\subseteq T$ be nonempty, and let $N(s)$ be the
set of paths of $D$ with an end in $s$. The paths ending in $s$ carry at least
$\sum_{v\in s}b_D(v)\ge 2|s|-|s\cap O|$ ends in $s$, and each path has at most two ends, so
$2|N(s)|\ge 2|s|-|s\cap O|\ge 2|s|-|O|$. Hence $|N(s)|\ge|s|-|O|/2$, and since $|N(s)|$ is an
integer, $|N(s)|\ge|s|-k$. So $s$ has at least $|N(s)|+k\ge|s|$ neighbours. By Hall's theorem~\cite{Hall35} there
is an injective map $f$ from $T$ to the slots with $f(v)$ adjacent to $v$. Assign each path $P$ of
the form $P=f(v)$ to $v$. This is a marked partition of $(F,T)$: $v$ is an end of $f(v)$, and
distinct marks receive distinct paths because $f$ is injective. At most $k$ marks are sent to dummy
slots, so at least $|T|-k$ paths are assigned, and the cost is at most $|D|-|T|+k$.
\end{proof}

In applications $O$ is a set of odd vertices, which always have an end by parity, and the marks outside $O$ are even vertices with at least one end, hence with two.

\subsection{The marked tests}

\begin{definition}[tests and the bounded marked bound]\label{def:MB}
For $m\ge0$ we consider three assertions.
\begin{itemize}[leftmargin=4em]
\item[$(\TA_m)$] For every $2$-connected graph $F$ with $|V(F)|\le m$ and every nonempty
$S\subseteq V(F)$, the marked bound \eqref{eq:mb} holds.
\item[$(\TB_m)$] For every $2$-connected graph $J$ with $|V(J)|=2k+1\le m$ and $p(J)=k+1$, and for
every odd vertex $x$ of $J$, some decomposition of $J$ has exactly $k+1$ paths and at least three ends
at $x$.
\item[$(\MBU_m)$] For every connected graph $F$ whose $2$-connected blocks have order at most $m$,
and every nonempty $S\subseteq V(F)$, the marked bound \eqref{eq:mb} holds.
\end{itemize}
\leanst{Gallai.MarkedTwoConnTestUpTo, Gallai.ThreeEndsTestUpTo, Gallai.MarkedBoundUpTo}
\end{definition}

\begin{theorem}[block transfer; MR Theorem 1, finite form]\label{thm:transfer}
For every $m\ge0$, $(\TA_m)$ and $(\TB_m)$ imply $(\MBU_m)$.
\leanst{Gallai.mbOn\_of\_testsUpTo, Gallai.markedBoundUpTo\_of\_tests}
\end{theorem}

\begin{remark}
With $m$ unbounded, Theorem~\ref{thm:transfer} and Lemma~\ref{lem:marked-basic}(a) show that the
two tests on all $2$-connected graphs imply Gallai's conjecture. MR Theorem~1 also proves the
converse, so that Gallai's conjecture is equivalent to the two tests (Lean:
\texttt{Gallai.gallai\_iff\_tests}). The spine needs only the sufficiency, and only in the
finite-order form above, because Theorem~59 feeds it tests of bounded order.
\end{remark}

\section{Block transfer and matching cuts}\label{sec:transfer}

This section proves Theorem~\ref{thm:transfer} (MR~\S2, ``Bridges, pendant edges and cut
vertices'' and ``Transfer through a block tree'') and the matching-cut case of MR Theorem~7 with
Lemma~6 (Theorem~\ref{thm:mc}).

\subsection{Partial marked partitions}
Let $G$ be a graph, $E'\subseteq E(G)$ and $T\subseteq V(G)$. A \emph{partial marked partition of
$E'$ with marks $T$} is a list $M$ of paths of $G$ whose edge sets partition $E'$, together with an
injective partial assignment $\sigma$ of paths to marks in $T$, each assigned mark being an end of
its path; its cost $c(M)$ is the number of unassigned paths \lean{Gallai.PMP}. If $U\subseteq V(G)$
and $E'=E(G[U])$, then every path of $M$ lies in $G[U]$ (each of its vertices is on one of its
edges), so a partial marked partition of $E(G[U])$ with marks $T\subseteq U$ is exactly a marked
partition of $(G[U],T)$. We say a mark $x$ is \emph{used} by $M$ if some path is assigned to it.

\begin{lemma}[gluing at a vertex]\label{lem:glue}
Let $U_1,U_2\subseteq V(G)$ with $U_1\cap U_2=\{x\}$, and let $E_1,E_2$ be disjoint edge sets with
both ends of every edge of $E_i$ in $U_i$. Let $M_i$ be a partial marked partition of $E_i$ with
marks $T$ and cost $c_i$. Then $E_1\cup E_2$ has a partial marked partition $M$ as follows.
\begin{enumerate}[label=\textup{(G\arabic*)},leftmargin=3.2em]
\item If $x$ is used by at most one of $M_1,M_2$: marks $T$ and $c(M)=c_1+c_2$.
\item In general: marks $T$ and $c(M)\le c_1+c_2+1$.
\item If $M_1$ has a path $P$ ending at $x$ that is not assigned to $x$: marks $T$ and
$c(M)\le c_1+c_2$.
\item In general: marks $T\setminus\{x\}$ and $c(M)\le c_1+c_2+1$.
\item If $x$ is used by at most one of $M_1,M_2$, $M_1$ has an unassigned path $P$ ending at $x$, and
$M_2$ has a path $R$ ending at $x$ that is not assigned to $x$: marks $T$ and $c(M)\le c_1+c_2-1$.
\item If $x$ is not used by $M_2$, $M_1$ has two distinct unassigned paths $P,Q$ ending at $x$, and
$M_2$ has a path $R$ with $x$ as an internal vertex: marks $T$ and $c(M)\le c_1+c_2-1$.
\end{enumerate}
\leanst{Gallai.glue0, Gallai.glue1, Gallai.glueJ, Gallai.glueU, Gallai.glueEnd,
Gallai.glueSplit}
\end{lemma}

\begin{proof}
Every path of $M_i$ has all its vertices in $U_i$, so a path of $M_1$ and a path of $M_2$ meet at
most in $x$, and a mark used by both $M_1$ and $M_2$ lies in $U_1\cap U_2$, hence equals $x$.

(G1) Take the union; the assignment stays injective since $x$ is the only possible common mark.
(G2) Unassign the path of $M_1$ assigned to $x$ (if any), at cost at most one, and apply (G1).
(G3) If $x$ is not used by $M_2$, apply (G1). Otherwise let $Q\in M_2$ be assigned to $x$; then $Q$
ends at $x$. Replace $P$ and $Q$ by the path $P\cup Q$ (Lemma~\ref{lem:split}(c)) and give it the
assignment of $P$. That assignment, if any, is an end of $P$ other than $x$, hence an end of $P\cup
Q$. The path $Q$ was assigned, so the cost is $c_1+c_2$; the mark $x$ is now used at most by $M_1$.
(G4) Unassign $x$ on both sides. If $x$ was used by both, by $P_1\in M_1$ and $P_2\in M_2$, replace
them by the unassigned path $P_1\cup P_2$; two assigned paths become one unassigned path. Otherwise
at most one path becomes unassigned.
(G5) Take the union (G1), and replace $P$ and $R$ by $P\cup R$ with the assignment of $R$, which is
an end of $R$ other than $x$. The unassigned path $P$ disappears and the assignment status of $R$ is
inherited, so the cost drops by one.
(G6) Take the union (G1). Split $R$ at $x$ into $R_1,R_2$, each ending at $x$; the assignment of $R$,
if any, is an end of $R$ other than $x$ and goes to the half containing it. Replace $P,Q,R$ by
$P\cup R_1$ and $Q\cup R_2$, where $P\cup R_1$ inherits the assignment of $R_1$ and $Q\cup R_2$ that
of $R_2$. Before, $P$ and $Q$ cost two and $R$ cost $\ind{R\text{ unassigned}}$; after, the two new
paths cost $1+\ind{R\text{ unassigned}}$.
\end{proof}

\begin{lemma}[pendant edges]\label{lem:pendant-edge}
Let $M$ be a partial marked partition of $E'$ with marks $T$ and cost $c$.
\begin{enumerate}[label=\textup{(\alph*)}]
\item \emph{(trimming)} Let $xy\in E'$ be the only edge of $E'$ at $y$, with $y\notin T$. Then
$E'\setminus\{xy\}$ has a partial marked partition $M'$ with marks $T$ such that one of the following
holds: (i) $c(M')=c-1$, and $M'$ has one path fewer than $M$; (ii) $c(M')=c$, $M'$ has one path
fewer than $M$, and $x$ is not used by $M'$; (iii) $c(M')=c$, $M'$ has as many paths as $M$, and
$M'$ has a path ending at $x$ that is not assigned to $x$.
\item \emph{(attaching)} Let $wv\in E(G)\setminus E'$ where $w$ lies on no path of $M$, and let
$o$ be either ``none'' or the mark $w$, where in the second case $w\in T\setminus\{v\}$ and $w$ is not
used by $M$. Then $E'\cup\{wv\}$ has a partial marked partition with marks $T\setminus\{v\}$ and cost
$c+\ind{o=\text{none}}$.
\end{enumerate}
\leanst{Gallai.PMP.exists\_trim, Gallai.attach\_pendant}
\end{lemma}

\begin{proof}
(a) Let $Q$ be the path containing $xy$; since $y$ has no other edge, $y$ is an end of $Q$. If
$Q=xy$, then $Q$ is unassigned or assigned to $x$ (as $y\notin T$); deleting $Q$ gives (i) or (ii).
Otherwise $Q=yx\cdots t$ with $t\ne x$, and $Q$ is unassigned or assigned to $t$. Deleting $y$ from
$Q$ gives a path ending at $x$ with the same assignment, which is (iii).
(b) If $v$ is used, by $P$, then $P$ ends at $v$ and avoids $w$; replace $P$ by $P\cup vw$ with
assignment $o$. Otherwise add the path $wv$ with assignment $o$. In both cases $v$ is no longer used,
the only assignment introduced is $o$, and the cost changes by $\ind{o=\text{none}}$ (in the first
case $P$ was assigned).
\end{proof}

\subsection{Proof of Theorem~\ref{thm:transfer}}
Fix $m$ and assume $(\TA_m)$ and $(\TB_m)$. We prove by strong induction on $n$:
\begin{quote}
$(\ast)_n$: if $G$ is connected of order $n$ and its $2$-connected blocks have order at most $m$,
then $q(G,S)\le\fl{(n-|S|)/2}$ for every nonempty $S\subseteq V(G)$.
\end{quote}
A $2$-connected induced subgraph of an induced subgraph of $G$ is a $2$-connected induced subgraph of
$G$. So the induction hypothesis gives the marked bound for every connected induced subgraph
$G[U']$ with $|U'|<n$ and every nonempty set of marks in $U'$. Let $G$ and $S$ be as in $(\ast)_n$.
We go through the steps below in order; each step either proves the bound or passes the
information that it did not apply to the next one \lean{Gallai.mbOn\_of\_testsUpTo}.

\emph{Step 0: $n\le2$} \lean{Gallai.baseStep}. Then $G$ has no edge or a single edge $ab$ with
$\{a,b\}=V(G)$. In the second case assign the path $ab$ to a mark. Either way $q(G,S)=0$.

\emph{Step 1: $G$ is $2$-connected} \lean{Gallai.twoConnStep}. Then $n\le m$ and $(\TA_m)$
applies.

\emph{Step 2: $G$ has a cut vertex $x$ with a separation $(C,D)$ such that $|C|,|D|\ge2$.} This is
Lemma~\ref{lem:cut} below.

\emph{Step 3: $G$ has a leaf $w$, with neighbour $v$, and it is not the case that both $S=\{v\}$
and $n$ is even.} This is Lemma~\ref{lem:leafstep}.

\emph{Step 4: $S=\{v\}$ and $G$ has two distinct leaves $w,w'$ adjacent to $v$.} This is
Lemma~\ref{lem:doubleleaf}.

\emph{Step 5: none of Steps 0--4 applies.} Then $n\ge3$ and $G$ is not $2$-connected, so $G$ has a
cut vertex $x$ and a separation $(C,D)$ at $x$. Since Step~2 does not apply, one side, say $C=\{a\}$,
is a single vertex; then $N_G(a)\subseteq\{x\}$, and $a$ is a leaf at $x$ since $G$ is connected. So
$G$ has a leaf. By Step~3, every leaf $w$ is adjacent to the unique mark $v$, where $S=\{v\}$ and $n$
is even. By Step~4 there is only one leaf, $w$. Now $J=G-w$ is connected with $n-1\ge3$ vertices
(as $n$ is even and at least three). If $J$ had a cut vertex, Lemma~\ref{lem:residue} would give a
separation of $G$ as in Step~2 or a second leaf of $G$; both are excluded. So $J$ is $2$-connected,
and Lemma~\ref{lem:pendant} gives the bound. This completes the induction, given the five lemmas
below. Theorem~\ref{thm:transfer} is $(\ast)_n$ for all $n$.\qed

\medskip
In Lemmas~\ref{lem:cut}--\ref{lem:pendant}, $G$ and $S$ are as in $(\ast)_n$ and the marked bound
holds for every connected induced subgraph of order less than $n$. Lemma~\ref{lem:residue} is a
statement about any connected graph.

\begin{lemma}[cut vertices; MR p.~5]\label{lem:cut}
If $G$ has a separation $(C,D)$ at $x$ with $|C|,|D|\ge2$, then $q(G,S)\le\fl{(n-|S|)/2}$.
\leanst{Gallai.cutStep}
\end{lemma}

\begin{proof}
Put $U_1=C\cup\{x\}$, $U_2=D\cup\{x\}$ and $G_i=G[U_i]$. Then $|U_i|\ge3$, $U_1\cap U_2=\{x\}$,
$|U_1|+|U_2|=n+1$, and $E(G)=E(G_1)\sqcup E(G_2)$. Each $G_i$ is connected: for $c\in C$, a path of $G$
from $c$ to a vertex of $D$ must pass through $x$, and its part from $c$ to the first visit of $x$ lies
in $G_1$; so every vertex of $C$ is joined to $x$ in $G_1$, and symmetrically every vertex of $D$ is
joined to $x$ in $G_2$. In
particular $x$ has a neighbour $y_i\in U_i\setminus\{x\}$. Let $S_i=S\cap U_i$. The graph
$G_1^+=G[U_1\cup\{y_2\}]$ is $G_1$ plus the pendant edge $xy_2$ (the only neighbour of $y_2$ in
$U_1$ is $x$). It is connected and has $|U_1|+1\le n-1$ vertices, since $|U_2|\ge3$. The same holds
for $G_2^+=G[U_2\cup\{y_1\}]$. All the graphs $G_1,G_2,G_1^+,G_2^+$ are connected induced subgraphs of
order less than $n$. We glue with Lemma~\ref{lem:glue} (with $E_i=E(G_i)$ and the marks $T$ named in
each case), and the result is a marked partition of $(G,S)$ whenever its marks are $S$.

\emph{Case 1: $x\in S$.} Let $r_i=|U_i|-|S_i|$, the number of unmarked vertices of $U_i$. Then
$n-|S|=r_1+r_2$ because $x$ is counted in both $S_1$ and $S_2$.

(1a) Both $r_i$ odd. Marked partitions $M_i$ of $(G_i,S_i)$ cost at most $(r_i-1)/2$. By (G2) the
union costs at most $(r_1-1)/2+(r_2-1)/2+1=(r_1+r_2)/2$.

(1b) $r_1$ even (the case $r_2$ even is symmetric). The mark set $S_1$ does not contain $y_2$, and
$q(G_1^+,S_1)\le\fl{(r_1+1)/2}=r_1/2$. Take an optimal $M^+$ and trim the pendant edge $xy_2$
(Lemma~\ref{lem:pendant-edge}(a), $y_2\notin S_1$). Let $M_2$ be optimal for $(G_2,S_2)$, of cost at
most $\fl{r_2/2}$. In outcome (i) the trimmed $M_1$ costs at most $r_1/2-1$, and (G2) gives at most
$r_1/2+\fl{r_2/2}$. In outcome (ii) $x$ is not used by $M_1$, and (G1) gives the same bound. In
outcome (iii) $M_1$ has a path ending at $x$ not assigned to $x$, and (G3) gives the same bound. As
$r_1$ is even, $r_1/2+\fl{r_2/2}=\fl{(r_1+r_2)/2}$.

\emph{Case 2: $x\notin S$.} Let $r_i=|U_i|-1-|S_i|$, the number of unmarked vertices of $U_i$ other
than $x$. Then $n-|S|=r_1+r_2+1$.

(2a) Some $r_i$ odd, say $r_1$. Mark $x$ on both sides: $q(G_i,S_i\cup\{x\})\le\fl{r_i/2}$. By (G4)
with $T=S\cup\{x\}$ the union is a marked partition of $(G,S)$ of cost at most
$\fl{r_1/2}+\fl{r_2/2}+1=(r_1+1)/2+\fl{r_2/2}=\fl{(r_1+r_2+1)/2}$.

(2b) Both $r_i$ even and both $S_i$ nonempty. Then $q(G_i,S_i)\le\fl{(r_i+1)/2}=r_i/2$, $x$ is used
by neither side, and (G1) gives $(r_1+r_2)/2=\fl{(r_1+r_2+1)/2}$.

(2c) Both $r_i$ even and $S_1=\emptyset$ (the case $S_2=\emptyset$ is symmetric). Then $S=S_2\ne\emptyset$,
$r_1=|U_1|-1$, and $M_2$ optimal for $(G_2,S_2)$ costs at most $r_2/2$. In $G_1^+$ mark only $x$:
$q(G_1^+,\{x\})\le\fl{|U_1|/2}=r_1/2$. Take an optimal $M^+$ (marks $\{x\}$, at most one assigned
path) and trim $xy_2$. We claim that $G_1$ has a decomposition $M_1$, viewed as a partial marked
partition with marks $S$ and no assigned path, such that either
($\alpha$) $|M_1|\le c(M^+)$, or
($\beta$) $|M_1|\le c(M^+)+1$ and $M_1$ has two distinct paths ending at $x$.
In outcome (i) a path was deleted and $|M^+|\le c(M^+)+1$, so ($\alpha$) holds. In outcome (ii) the
deleted path was the only assigned one, so $|M_1|=|M^+|-1=c(M^+)$. In outcome (iii) the path $L$
ending at $x$ is unassigned (its assignment is not $x$, and $x$ is the only mark). If $x$ is used by
some path $Q\ne L$ of the trimmed partition, then $L$ and $Q$ both end at $x$ and
$|M_1|=|M^+|=c(M^+)+1$, which is ($\beta$). Otherwise nothing is assigned and $|M_1|=c(M^+)$.

In case ($\alpha$), (G1) gives cost at most $r_1/2+r_2/2$. In case ($\beta$), $M_1$ has two
unassigned paths $P,Q$ ending at $x$, and $x$ is used by neither side ($x\notin S$). If some path $R$
of $M_2$ ends at $x$, (G5) gives cost at most $(r_1/2+1)+r_2/2-1$. Otherwise no path of $M_2$ ends at
$x$, while $x$ has the edge $xy_2$ in $G_2$; by Lemma~\ref{lem:split}(b) some path of $M_2$ has $x$
as an internal vertex, and (G6) gives the same bound. In all cases the cost is at most
$(r_1+r_2)/2=\fl{(n-|S|)/2}$.

The cases are exhaustive: for $x\in S$, either some $r_i$ is even or both are odd; for $x\notin S$,
either some $r_i$ is odd, or both are even and then both $S_i$ are nonempty or one is empty (not
both, as $S\ne\emptyset$).
\end{proof}

\begin{lemma}[leaves; MR pp.~4--5]\label{lem:leafstep}
Let $n\ge3$ and let $w$ be a leaf of $G$ with neighbour $v$. Unless $S=\{v\}$ and $n$ is even,
$q(G,S)\le\fl{(n-|S|)/2}$.
\leanst{Gallai.leafStep}
\end{lemma}

\begin{proof}
$G-w$ is connected of order $n-1$ (a path between two vertices other than $w$ cannot pass through
the leaf $w$), and $E(G)=E(G-w)\sqcup\{wv\}$. In each case we take an optimal
marked partition $M$ of $G-w$ for the stated marks and add the edge $wv$.
\begin{enumerate}[label=(L\arabic*),leftmargin=3em]
\item $w\in S$ and $S\ne\{w\}$. Marks $S\setminus\{w\}$: $c(M)\le\fl{(n-1-(|S|-1))/2}$. Add the path
$wv$ assigned to $w$.
\item $S=\{w\}$. Marks $\{v\}$: $c(M)\le\fl{(n-2)/2}$. Attach $wv$ with $o=w$
(Lemma~\ref{lem:pendant-edge}(b), marks $\{v,w\}$); the result has marks $\{w\}$ and the same cost,
at most $\fl{(n-1)/2}$.
\item $w\notin S$, $v\in S$ and $S\ne\{v\}$. Marks $S\setminus\{v\}$:
$c(M)\le\fl{(n-1-(|S|-1))/2}$. Add the path $vw$ assigned to $v$.
\item $S=\{v\}$; then $n$ is odd by hypothesis. Marks $\{v\}$: $c(M)\le\fl{(n-2)/2}=(n-3)/2$.
Attach $wv$ with $o=$ none: cost at most $(n-1)/2=\fl{(n-1)/2}$.
\item $w,v\notin S$. Marks $S\cup\{v\}$: $c(M)\le\fl{(n-1-|S|-1)/2}=\fl{(n-|S|)/2}-1$. Attach $wv$
with $o=$ none; the result has marks $S$ and cost at most $\fl{(n-|S|)/2}$.
\end{enumerate}
These cases are exhaustive.
\end{proof}

\begin{lemma}[two leaves at the only mark]\label{lem:doubleleaf}
If $S=\{v\}$ and $w\ne w'$ are leaves adjacent to $v$, then $q(G,\{v\})\le\fl{(n-1)/2}$.
\leanst{Gallai.doubleLeafStep}
\end{lemma}

\begin{proof}
$G-w-w'$ is connected (no path between two other vertices passes through a leaf), contains $v$, and
has order $n-2$. An optimal marked partition of
$(G-w-w',\{v\})$ costs at most $\fl{(n-3)/2}$. Add the unassigned path $wvw'$.
\end{proof}

\begin{lemma}[the pendant configuration; MR p.~5]\label{lem:pendant}
Let $n$ be even, $S=\{v\}$, $w$ a leaf at $v$, and suppose $J=G-w$ is $2$-connected. Then
$q(G,\{v\})\le\fl{(n-1)/2}$.
\leanst{Gallai.pendantStep}
\end{lemma}

\begin{proof}
Write $n=2k+2$, so $|V(J)|=2k+1$, and the target is $k$. As $J$ is a $2$-connected induced subgraph
of $G$, whose $2$-connected blocks have order at most $m$, we have $2k+1\le m$. By $(\TA_m)$,
$q(J,\{v\})\le\fl{2k/2}=k$, so $p(J)\le k+1$ by Lemma~\ref{lem:marked-basic}(a). If $p(J)\le k$,
take an optimal decomposition of $J$ and add the path $vw$ assigned to $v$; the cost is at most $k$.
Otherwise $p(J)=k+1$, and we find a decomposition $D$ of $J$ with $|D|=k+1$ and $b_D(v)\ge2$. If
$d_J(v)$ is odd this is $(\TB_m)$. If $d_J(v)$ is even, take an optimal marked partition
$(D,\sigma)$ of $(J,\{v\})$; it costs at most $k$ and has at most one assigned path, while
$|D|\ge p(J)=k+1$. Hence $|D|=k+1$ and $v$ is used, so $b_D(v)\ge1$, and $b_D(v)\ge2$ by parity.
Now two distinct paths $P_1,P_2$ of $D$ end at $v$. Assign $P_2$ to $v$ and extend $P_1$ by the edge
$vw$. This is a marked partition of $(G,\{v\})$ with $k+1$ paths, one assigned, so it costs $k$.
\end{proof}

As in MR, an even $d_J(v)$ causes no difficulty: an optimal marked partition already has two ends
at $v$.

\begin{lemma}[the residue; MR p.~6]\label{lem:residue}
Let $w$ be a leaf of a connected graph $G$ with neighbour $v$, and suppose $J=G-w$ has a cut vertex
$x'$ with a separation $(C',D')$. Then $G$ has a separation with both sides of size at least two, or
$G$ has a leaf other than $w$.
\leanst{Gallai.residue\_sep}
\end{lemma}

\begin{proof}
If $x'=v$: then $(C',D'\cup\{w\})$ is a separation of $G$ at $v$, since $w$ is adjacent only to $v$.
The second side has at least two vertices. If $|C'|\ge2$ we are done; otherwise $C'=\{a\}$, and $a$
is a leaf of $G$ at $v$ with $a\ne w$. If $x'\ne v$: by symmetry $v\in C'$, and $(C'\cup\{w\},D')$ is
a separation of $G$ at $x'$, since the only neighbour of $w$ lies in $C'$. The first side has at
least two vertices. If $|D'|\ge2$ we are done; otherwise $D'=\{b\}$ is a leaf of $G$ at $x'$, and
$b\ne w$. (A singleton side $\{a\}$ of a separation at $x$ in a connected graph is a leaf at $x$,
since all its neighbours lie in $\{x\}$.)
\end{proof}

\subsection{Matching cuts}

\begin{lemma}[connected shores]\label{lem:shores}
A connected graph $B$ with a matching cut has a matching cut $(X,Y)$ such that $B[X]$ and $B[Y]$ are
connected.
\leanst{Gallai.exists\_connected\_matchingCut}
\end{lemma}

\begin{proof}
Let $(X_0,Y_0)$ be a matching cut. Say a set $Z$ has the \emph{matching property} if every vertex has
at most one neighbour on the other side of $(Z,V(B)\setminus Z)$; this is symmetric in $Z$ and its
complement. If $Z$ has it and $K\subseteq Z$ is a union of components of $B[Z]$, then $K$ has it: a
vertex of $K$ has no neighbour in $Z\setminus K$, so its neighbours outside $K$ are its neighbours
outside $Z$; a vertex of $Z\setminus K$ has no neighbour in $K$; and a vertex outside $Z$ has its
neighbours in $K$ among its neighbours in $Z$. Let $C$ be the component of $B[X_0]$ containing some
$y_0\in X_0$; then $C$ has the matching property. Pick $z_0\in Y_0$, so $z_0\notin C$, and let $X$ be
the component of $B-C$ containing $z_0$; it has the matching property because $V(B)\setminus C$
does. Put $Y=V(B)\setminus X$. Then $X\ni z_0$ and $Y\supseteq C\ni y_0$ are nonempty, $B[X]$ is
connected, and $B[Y]$ is connected: it consists of $C$ and the components of $B-C$ other than $X$,
and each of these has an edge to $C$ because $B$ is connected and its only neighbours outside itself
lie in $C$.
\end{proof}

For $A\subseteq V(B)$ let $L_A(B)$ be $B$ with a private leaf $a'$ attached at every $a\in A$
\lean{Gallai.leafGraph}.

\begin{lemma}[trimming the leaves]\label{lem:trimleaves}
Every decomposition $D$ of $L_A(B)$ yields a decomposition $D'$ of $B$ with
$|D'|+|\{a\in A:b_{D'}(a)=0\}|\le|D|$.
\leanst{Gallai.leafGraph\_trim}
\end{lemma}

\begin{proof}
A leaf edge $aa'$ can occur in a path only at one of its ends, with $a'$ as the end. Delete all leaf
edges from every path. A path consisting only of leaf edges is a single leaf edge, and it disappears;
every other path $P$ becomes a path $P^\circ$ of $B$, and for each leaf edge $aa'$ removed from $P$,
the vertex $a$ is an end of $P^\circ$. The result $D'$ is a decomposition of $B$. If $a\in A$ has
$b_{D'}(a)=0$, then the path of $D$ containing $aa'$ cannot have survived, so it was the single edge
$aa'$. These deleted paths are distinct for distinct $a$, which gives the inequality.
\end{proof}

\begin{lemma}[MR Lemma 6, matching-cut form]\label{lem:lemma6}
Let $(X,Y)$ be a matching cut of $B$ with crossing edges $a_1b_1,\dots,a_tb_t$, where $a_i\in X$ and
$b_i\in Y$, and let $A\subseteq V(B)$ and $H=L_A(B)$. Let $H_X$ be the subgraph of $H$ induced on $X$
and the leaves of $A\cap X$, and $H_Y$ likewise. If $M_X$ is a marked partition of
$(H_X,\{a_1,\dots,a_t\})$ and $M_Y$ one of $(H_Y,\{b_1,\dots,b_t\})$, then
$p(H)\le c(M_X)+c(M_Y)+t$.
\leanst{Gallai.lemma6}
\end{lemma}

\begin{proof}
The $a_i$ are distinct and the $b_i$ are distinct, by the matching property. The edges of $H$ are
those of $H_X$, those of $H_Y$, and the $t$ crossing edges. For each $i$, let $P_i$ be the path of
$M_X$ assigned to $a_i$ and $Q_i$ the path of $M_Y$ assigned to $b_i$, when they exist, and replace
them by the single path $P_i\cdot a_ib_i\cdot Q_i$ (just $a_ib_i$, or $a_ib_i$ extended on one side,
when one or both do not exist). This is a path because $P_i\subseteq H_X$ and $Q_i\subseteq H_Y$ are
vertex-disjoint and end at $a_i$ and $b_i$. Distinct $i$ use distinct paths, since the assignments are
injective. The result is a decomposition of $H$ consisting of the $c(M_X)+c(M_Y)$ unassigned paths and
$t$ new paths.
\end{proof}

\begin{theorem}[MR Theorem 7, matching-cut case]\label{thm:mc}
Let $B$ be a connected graph of order $n$ with a matching cut, and assume $(\MBU_{n-1})$. Then
$q(B,S)\le\fl{(n-|S|)/2}$ for every $S\subseteq V(B)$, including $S=\emptyset$. In particular
$p(B)\le\fl{n/2}$.
\leanst{Gallai.matchingCut\_markedBound}
\end{theorem}

\begin{proof}
Take a matching cut $(X,Y)$ with connected shores (Lemma~\ref{lem:shores}), with crossing edges
$a_ib_i$ as in Lemma~\ref{lem:lemma6}; $t\ge1$ since $B$ is connected. Split $S=A\sqcup O$ into its
even vertices $A$ and odd vertices $O$; let $\alpha=|A|$ and $o=|O|$. Let $H=L_A(B)$ and $H_X,H_Y$ as in
Lemma~\ref{lem:lemma6}. The graph $H_X$ is connected ($B[X]$ is, and leaves are attached to it). A
$2$-connected induced subgraph of $H_X$ has minimum degree at least two, so it contains no leaf and
lies in $B[X]$; hence the $2$-connected blocks of $H_X$ have order at most $|X|\le n-1$. The same holds
for $H_Y$. By $(\MBU_{n-1})$ with the nonempty mark sets $\{a_i\}$ and $\{b_i\}$,
\[
q(H_X,\{a_i\})\le\fl{\tfrac{|X|+|A\cap X|-t}2},\qquad q(H_Y,\{b_i\})\le\fl{\tfrac{|Y|+|A\cap Y|-t}2}.
\]
By Lemma~\ref{lem:lemma6}, $H$ has a decomposition $D$ with
$|D|\le(n+\alpha-2t)/2+t=(n+\alpha)/2$, hence $|D|\le\fl{(n+\alpha)/2}$. By
Lemma~\ref{lem:trimleaves} there is a decomposition $D'$ of $B$ with $|D'|+\zeta\le|D|$, where
$\zeta=|\{v\in A:b_{D'}(v)=0\}|$. Let $T=\{v\in A:b_{D'}(v)\ge1\}\cup O$, so
$|T|=\alpha-\zeta+o$. Every vertex of $O$ has odd degree, hence an end, and every vertex of $T\setminus O$ is
an even vertex with an end, hence with two (Lemma~\ref{lem:parity}). Lemma~\ref{lem:hall} and
Lemma~\ref{lem:marked-basic}(b) give
\[
q(B,S)\le q(B,T)\le|D'|+\fl{o/2}-|T|\le\fl{\tfrac{n+\alpha}2}-\zeta+\fl{\tfrac o2}-(\alpha-\zeta+o)
=\fl{\tfrac{n-\alpha}2}-\cl{\tfrac o2}\le\fl{\tfrac{n-\alpha-o}2}.
\]
The last step holds because $\fl{\mu}+\cl{\nu}\ge\fl{\mu+\nu}$ with $\mu=(n-\alpha-o)/2$ and
$\nu=o/2$. For $S=\emptyset$ this reads $p(B)=q(B,\emptyset)\le\fl{n/2}$.
\end{proof}

\section{The class \texorpdfstring{$\cC$}{C} and Theorem 59}\label{sec:thm59}

\begin{definition}[rooted assertions; MR \S17]\label{def:rooted}
Let $B$ be a graph of order $n$ and $x\in V(B)$. We say $B$ \emph{passes the rooted test at $x$},
written $\RT(B,x)$, if both of the following hold:
\begin{enumerate}[label=\textup{(R\arabic*)},leftmargin=3em]
\item if $n$ and $d_B(x)$ are even, some decomposition $D$ of $B$ has $|D|\le n/2$ and $b_D(x)\ge2$;
\item if $n$ and $d_B(x)$ are odd, some decomposition $D$ of $B$ has $|D|\le\fl{n/2}+1$ and
$b_D(x)\ge3$.
\end{enumerate}
For $N\ge0$, $(\RT_N)$ is the assertion that $\RT(B,x)$ holds for every $B\in\cC$ with
$|V(B)|\le N$ and every $x\in V(B)$.
\leanst{Gallai.RootedAt, Gallai.RootedTestsUpTo}
\end{definition}

In MR notation, (R1) is $F_2(B,x)\le k$ for $n=2k$ and (R2) is $F_3(B,x)\le k+1$ for $n=2k+1$.

\begin{theorem}[MR Theorem 59, finite-order form]\label{thm:59}
For every $m\ge0$, $(\RT_{m+1})$ implies $(\MBU_m)$.
\leanst{Gallai.thm59 : Gallai.Thm59Statement}
\end{theorem}

This is the form recommended by audit A2 (finding F1): the conclusion is the marked bound
\eqref{eq:mb}, not only Gallai's bound, on every connected graph whose $2$-connected blocks have
order at most $m$. The proof uses the rooted tests only on members of $\cC$ of order $n\le m$ and on
cones $B+z(S)$ over $|S|\ge2$ marks of such members, of order $n+1\le m+1$ (A1, G10).

\begin{lemma}[cones stay in $\cC$]\label{lem:coneC}
If $B\in\cC$ and $S\subseteq V(B)$ with $|S|\ge2$, then $B+z(S)\in\cC$.
\leanst{Gallai.apex\_inClassC}
\end{lemma}

\begin{proof}
Let $X=B+z(S)$; it has at least four vertices. $X-z=B$ is connected; $X-v$ for $v\in V(B)$ consists
of the connected graph $B-v$ and the vertex $z$, which keeps a neighbour in $S\setminus\{v\}\ne\emptyset$;
and $X$ itself is connected. So $X$ is $2$-connected. Let $(P,Q)$ be a matching cut of $X$. If both
$P\cap V(B)$ and $Q\cap V(B)$ were nonempty, they would form a matching cut of $B$. So $V(B)$ lies
on one side, say $V(B)\subseteq P$, and then $Q=\{z\}$; but $z$ has $|S|\ge2$ neighbours in $P$, a
contradiction.
\end{proof}

\begin{proposition}[the marked tests on $\cC$; MR \S17, first part of the proof of Theorem 59]
\label{prop:classC}
Assume $(\RT_{m+1})$ and let $B\in\cC$ have order $n\le m$. Then:
\begin{enumerate}[label=\textup{(\alph*)}]
\item $q(B,S)\le\fl{(n-|S|)/2}$ for every nonempty $S\subseteq V(B)$;
\item if $n=2k+1$ and $p(B)=k+1$, then for every odd vertex $x$ some decomposition of $B$ has exactly
$k+1$ paths and at least three ends at $x$.
\end{enumerate}
\leanst{Gallai.markedBound\_of\_classC, Gallai.threeEnds\_of\_classC}
\end{proposition}

\begin{proof}
We prove (a) in three steps.

\emph{The cone step} \lean{Gallai.cone\_step}. Let $|S|=t\ge2$ with $n-t$ odd. By
Lemma~\ref{lem:coneC}, $X=B+z(S)\in\cC$, and $|V(X)|=n+1\le m+1$, $d_X(z)=t$. If $n$ is even, then
$t$ is odd and $n+1$ is odd, so (R2) for $X$ at $z$ gives a decomposition $D$ with
$|D|\le n/2+1$ and $b_D(z)\ge3$. Lemma~\ref{lem:cone} gives $t+2q(B,S)\le 2|D|-b_D(z)\le n-1$. If
$n$ is odd, then $t$ and $n+1$ are even, so (R1) gives $|D|\le(n+1)/2$ and $b_D(z)\ge2$, and again
$t+2q(B,S)\le n-1$. In both cases $q(B,S)\le(n-1-t)/2=\fl{(n-t)/2}$, since $n-t$ is odd.

\emph{Singletons} \lean{Gallai.singleton\_step}. Let $S=\{x\}$. If $n$ is even and $d_B(x)$ is
even, (R1) at $x$ gives $D$ with $|D|\le n/2$ and $b_D(x)\ge2$, so
$q(B,\{x\})\le n/2-1=\fl{(n-1)/2}$ by Lemma~\ref{lem:marked-basic}(d). If $n$ is even and $d_B(x)$ is
odd, we first find a decomposition $D$ with $|D|\le n/2$: by (R1) at any even vertex, or, if every
vertex is odd, by Lov\'asz's theorem (Theorem~\ref{thm:lovasz}). Then $b_D(x)\ge1$ by parity and
again $q(B,\{x\})\le n/2-1$. If $n$ is odd, pick $y\ne x$ (as $n\ge3$). Then $n-2$ is odd, the cone
step gives $q(B,\{x,y\})\le(n-3)/2$, and Lemma~\ref{lem:marked-basic}(c) gives
$q(B,\{x\})\le(n-1)/2$.

\emph{General $S$.} If $|S|=1$ use the singleton step, and if $|S|\ge2$ and $n-|S|$ is odd use the
cone step. If $|S|\ge2$ and $n-|S|$ is even, remove a mark $s$: $S'=S\setminus\{s\}$ has
$n-|S'|$ odd, and $q(B,S)\le q(B,S')$ by Lemma~\ref{lem:marked-basic}(b). If $|S'|\ge2$ the cone
step gives $q(B,S')\le(n-|S'|-1)/2=(n-|S|)/2$. If $|S'|=1$ then $n$ is even and the singleton step
gives $q(B,S')\le(n-2)/2=(n-|S|)/2$.

(b) By (R2) at $x$ some $D$ has $|D|\le k+1$ and $b_D(x)\ge3$, and $|D|\ge p(B)=k+1$.
\end{proof}

MR treats odd singletons by adjoining a vertex $w$ adjacent to $x$ and one other vertex $y$, and
applying (R1) at $w$; that graph is exactly the cone $B+z(\{x,y\})$, and the argument above replaces
the trimming by Lemma~\ref{lem:marked-basic}(c).

\begin{proof}[Proof of Theorem~\ref{thm:59}]
Assume $(\RT_{m+1})$. We prove $(\MBU_j)$ for all $j\le m$ by strong induction on $j$. Fix $j\le m$
and assume $(\MBU_{j'})$ for all $j'<j$. By Theorem~\ref{thm:transfer} it suffices to prove
$(\TA_j)$ and $(\TB_j)$. Let $F$ be $2$-connected of order $n\le j$.
\begin{itemize}
\item If $F$ has a matching cut, then $(\MBU_{n-1})$ holds by induction ($n-1<j$), and
Theorem~\ref{thm:mc} gives the marked bound for every $S$. For $S=\emptyset$ it gives
$p(F)\le\fl{n/2}$, so if $n=2k+1$ then $p(F)\le k$, and $F$ does not satisfy the hypothesis
$p(F)=k+1$ of $(\TB_j)$.
\item Otherwise $F\in\cC$ and $n\le m$, and Proposition~\ref{prop:classC} gives both assertions.
\end{itemize}
This proves $(\TA_j)$ and $(\TB_j)$, hence $(\MBU_j)$; for $j=m$ it is the theorem.
\end{proof}

The induction runs on the bound for the order of the $2$-connected blocks, not on the order of the
graph. This is necessary: the graphs $H_X,H_Y$ of Theorem~\ref{thm:mc} carry private leaves and may
have order at least $n$; only their blocks are smaller (audit~A14).

\begin{remark}[the converse]\label{rem:59conv}
Gallai's conjecture implies $\RT(B,x)$ for every $2$-connected $B$ and every $x$
(Lemma~\ref{lem:gallai-rooted} below, since $d_B(x)\ge2$ and an odd degree is then at least $3$).
Together with Theorem~\ref{thm:59} for all $m$ and a single mark, this is MR Theorem~59: Gallai's
conjecture is equivalent to $\RT(B,x)$ for all $B\in\cC$ and all $x$. We do not use the bridge and
diamond doubles by which MR expresses the two assertions as ordinary path-number bounds.
\end{remark}

\section{Separating pairs: MR Theorem 113}\label{sec:sep}

\begin{definition}\label{def:sepcomp}
Let $B$ be a graph and $x\ne y$ vertices. A set $C\subseteq V(B)\setminus\{x,y\}$ is a
\emph{separated component} for $\{x,y\}$ if $C$ is nonempty, $B[C]$ is connected, every neighbour of
a vertex of $C$ lies in $C\cup\{x,y\}$, and $V(B)\setminus(C\cup\{x,y\})\ne\emptyset$. Thus $C$ is
the vertex set of a component of $B-\{x,y\}$ and $\{x,y\}$ separates $C$ from the rest.
\leanst{Gallai.IsSepComponent}
\end{definition}

\begin{theorem}[separator theorem; MR Theorem 113, a strengthening of A2, F7]\label{thm:sep}
Let $B$ be a $2$-connected graph of order $n$ and assume $(\MBU_{n-2})$. Let $C$ be a separated
component for $\{x,y\}$ all of whose vertices are odd in $B$, and let $r\in V(B)\setminus C$ (possibly
$r\in\{x,y\}$). Then $\RT(B,r)$ holds.
\leanst{Gallai.separatorStatement : Gallai.SeparatorStatement, Gallai.separator\_theorem}
\end{theorem}

MR states this at a smallest failure of the rooted tests; its proof uses only that $B$ is
$2$-connected and $(\MBU_{n-2})$, which Theorem~\ref{thm:59} provides at a smallest failure. Neither
minimality nor the absence of a matching cut is assumed.

Throughout this section $B,n,x,y,C,r$ are as in Theorem~\ref{thm:sep}, and $V=V(B)$.

\subsection{Setup}
Let $A$ be the spanning subgraph of $B$ whose edges are the edges with at least one end in $C$,
together with $xy$ if $xy\in E(B)$. Let $E$ be the spanning subgraph whose edges are the remaining
edges of $B$: those with no end in $C$, other than $xy$ \lean{Gallai.sideA, Gallai.sideE}. Thus
$E(B)=E(A)\sqcup E(E)$ and $d_B=d_A+d_E$. All edges of $A$ lie inside $C\cup\{x,y\}$; $d_A(c)=d_B(c)$
and $d_E(c)=0$ for $c\in C$; and $d_A(v)=0$ for $v\notin C\cup\{x,y\}$. We call $x,y$ the
\emph{ports}.

\begin{lemma}\label{lem:sides}
\begin{enumerate}[label=\textup{(\alph*)}]
\item $E[V\setminus C]$ is connected.
\item $d_E(x),d_E(y),d_A(x),d_A(y)$ are all positive.
\item $|C|+\ind{d_A(x)\text{ odd}}+\ind{d_A(y)\text{ odd}}$ is even.
\item $|C|\ge2$.
\end{enumerate}
\leanst{Gallai.sideE\_induce\_connected, Gallai.sideE\_port\_pos, Gallai.sideA\_port\_pos,
Gallai.sideA\_parity, Gallai.two\_le\_card\_of\_odd}
\end{lemma}

\begin{proof}
(a) We first claim: if $\{p,q\}=\{x,y\}$, $v\notin C$ and $v\ne q$, then $v$ is joined to $p$ by a
path of $E$. Since $B$ is $2$-connected, $B-q$ has a path from $v$ to $p$. Follow it from $v$ until
it first reaches $p$. Before that it cannot enter $C$, because the only vertices outside $C$ with
neighbours in $C$ are $x$ and $y$, and $q$ is avoided. So this initial segment uses only edges with
no end in $C$, and it does not use $xy$ because it avoids $q$; it is a path of $E$. Now fix
$w\in V\setminus(C\cup\{x,y\})$. The claim with $p=x$, $q=y$ joins every $v\notin C$ with $v\ne y$,
and in particular $w$, to $x$; so all these vertices are joined to $w$. The claim with $p=y$, $q=x$
joins $w$ to $y$. Hence every vertex of $V\setminus C$ is joined to $w$ in $E$.

(b) By (a), $x$ and $y$ are joined in $E$ to a vertex $w\in V\setminus(C\cup\{x,y\})$, so they have positive degree in $E$. Pick
$c\in C$. In $B-y$ there is a path from $c$ to $w$; it leaves $C$ along an edge $c'u$ with $c'\in C$,
$u\notin C$, so $u\in\{x,y\}$ and $u\ne y$. Thus $c'x\in E(A)$ and $d_A(x)>0$; similarly $d_A(y)>0$.

(c) The odd vertices of $A$ are the vertices of $C$ (each has its odd $B$-degree) and those ports
that are odd in $A$; their number is even.

(d) If $C=\{c\}$, then $N_B(c)\subseteq\{x,y\}$ and $d_B(c)\ge2$ (Lemma~\ref{lem:kconn}), so
$d_B(c)=2$, which is even.
\end{proof}

By (c), if $|C|$ is odd exactly one port is odd in $A$, and if $|C|$ is even the ports have the same
parity in $A$. Write $|C|=2s+1$ or $|C|=2s$; by (d), $s\ge1$ in both cases.

\subsection{Smaller rooted bounds and the root analysis}

\begin{lemma}[smaller rooted bounds]\label{lem:small}
Assume $(\MBU_{n-2})$. Let $W$ be a graph with all edges inside $U\subseteq V(W)$, such that $W[U]$
is connected and $N:=|U|\le n-2$, and let $r\in U$.
\begin{enumerate}[label=\textup{(\alph*)}]
\item If $N$ is even and $d_W(r)$ is even and positive, some decomposition $D$ of $W$ has
$|D|\le N/2$ and $b_D(r)\ge2$.
\item If $N$ is odd and $d_W(r)$ is odd and at least $3$, some $D$ has $|D|\le\fl{N/2}+1$ and
$b_D(r)\ge3$.
\item If $N$ is odd and $d_W(r)=1$, then $p(W)\le\fl{N/2}$.
\end{enumerate}
\leanst{Gallai.R2\_on (part (a); the 2 counts ends), Gallai.R3\_on (part (b)), Gallai.floor\_on (part (c))}
\end{lemma}

\begin{proof}
We work in $W[U]$, which has the same edges and the same degrees on $U$. Its $2$-connected induced
subgraphs have at most $N\le n-2$ vertices, so $(\MBU_{n-2})$ applies to it.

(a) $q(W[U],\{r\})\le\fl{(N-1)/2}=N/2-1$. By Lemma~\ref{lem:marked-basic}(e) some $D$ has
$\lambda_D(r)\le N/2$, and Lemma~\ref{lem:ends}(a) finishes.

(b) Let $W'=W[U]^r$, with private leaf $\ell$. It is connected of order $N+1$, and its $2$-connected
induced subgraphs avoid $\ell$ (which has degree one), so they have at most $N\le n-2$ vertices.
Hence $q(W',\{r\})\le\fl{N/2}$, and by Lemma~\ref{lem:marked-basic}(e) some decomposition $D_1$ of
$W'$ has $\lambda_{D_1}(r)\le\fl{N/2}+1$. In $W'$ the vertex $r$ has even positive degree
$d_W(r)+1$, so Lemma~\ref{lem:ends}(a) gives $D_2$ with $|D_2|\le\fl{N/2}+1$ and $b_{D_2}(r)\ge2$.
Trim the leaf (Lemma~\ref{lem:leaf}(a)). Either the result $D$ has $|D|\le\fl{N/2}+1$ and
$b_D(r)\ge3$, or it has $|D|\le\fl{N/2}$ and $b_D(r)\ge1$; in the latter case
Lemma~\ref{lem:ends}(b) with $k=\fl{N/2}$ gives the claim.

(c) Let $a$ be the neighbour of $r$ and $T=U\setminus\{r\}$. Then $W[T]$ is connected (a path between
two vertices of $T$ cannot pass through the degree-one vertex $r$), $a\in T$, and $|T|=N-1$. Hence
$q(W[T],\{a\})\le\fl{(N-2)/2}=(N-3)/2$, and by Lemma~\ref{lem:marked-basic}(e) some decomposition
$D_0$ of $W[T]$ has $\lambda_{D_0}(a)\le(N-1)/2$. By Lemma~\ref{lem:leaf}(b),
$p(W[T]^a)\le(N-1)/2$. But $W[U]$ is $W[T]^a$ with $r$ as the private leaf, so
$p(W)\le(N-1)/2=\fl{N/2}$.
\end{proof}

\begin{lemma}[root analysis]\label{lem:root}
Assume $(\MBU_{n-2})$. Let $s\ge1$ and let $W$ be a graph on $V$ with all edges inside a set $U$,
such that $W[U]$ is connected, $|U|=n-2s$ and $r\in U$. Suppose that
\begin{enumerate}[label=\textup{(\roman*)}]
\item $d_W(r)\equiv d_B(r)\pmod2$ and $d_W(r)>0$, and
\item every decomposition $D_W$ of $W$ \emph{lifts}: $B$ has a decomposition $D$ with
$|D|\le|D_W|+s$ and $b_D(r)=b_{D_W}(r)$.
\end{enumerate}
Then $\RT(B,r)$ holds.
\leanst{Gallai.rootedAt\_of\_lift}
\end{lemma}

\begin{proof}
We have $|U|\le n-2$ and $|U|\equiv n\pmod 2$. (R1): if $n$ and $d_B(r)$ are even, then $|U|$ and
$d_W(r)$ are even, $d_W(r)>0$, and Lemma~\ref{lem:small}(a) gives $D_W$ with $|D_W|\le|U|/2$ and two
ends at $r$; its lift has at most $|U|/2+s=n/2$ paths. (R2): let $n$ and $d_B(r)$ be odd, so $|U|$
and $d_W(r)$ are odd. If $d_W(r)\ge3$, Lemma~\ref{lem:small}(b) and the lift give at most
$\fl{|U|/2}+1+s=\fl{n/2}+1$ paths and three ends at $r$. If $d_W(r)=1$, Lemma~\ref{lem:small}(c)
gives $D_W$ with at most $\fl{|U|/2}$ paths, whose lift $D$ has at most $\fl{n/2}$ paths. Since $B$
is $2$-connected, $d_B(r)\ge2$, and $d_B(r)$ is odd, so $d_B(r)\ge3$; Lemma~\ref{lem:ends}(b) turns
$D$ into a decomposition with at most $\fl{n/2}+1$ paths and three ends at $r$.
\end{proof}

\subsection{Path surgery}
In the constructions below a \emph{path} is any sequence of at least two distinct vertices of $V$,
regarded up to reversal; a family of paths whose edge sets partition $E(B)$ is a decomposition of $B$
\lean{Gallai.TPath, Gallai.exists\_decomp\_of\_edgesM}. Edge sets are added as multisets, so an
identity $E(M)=E(R_0)+E(A)$ for a family $M$ asserts in particular that no edge is repeated. For
paths $P,Q$ and a vertex $p$, write $P\cdot Q$ for the concatenation when $P$ ends and $Q$ starts at
$p$ and $V(P)\cap V(Q)=\{p\}$ (a path by Lemma~\ref{lem:split}(c)); a trivial piece consisting of $p$
alone is allowed.

\begin{lemma}[surgery]\label{lem:surgery}
\begin{enumerate}[label=\textup{(\alph*)}]
\item \emph{(substitution)} If a path $R$ contains the edge $ab$ and $T$ is a path from $a$ to $b$
with $V(T)\cap V(R)\subseteq\{a,b\}$, then replacing $ab$ in $R$ by $T$ gives a path $R'$ with the
same ends as $R$ and $E(R')=E(R)-ab+E(T)$.
\item \emph{(cut and attach)} Let $R=R_a\cdot ab\cdot R_b$, where $R_a$ ends at $a$ and $R_b$ starts
at $b$, and let $X$ and $Y$ be paths starting at $a$ and $b$ with $V(X)\cap V(R)\subseteq\{a,b\}$ and
$V(Y)\cap V(R)\subseteq\{a,b\}$. Then $N_1=R_a\cdot X$ and $N_2=Y^{-1}\cdot R_b$ are paths with
$E(N_1)+E(N_2)+ab=E(R)+E(X)+E(Y)$ and $b_{N_1}+b_{N_2}+\ind{\cdot=a}+\ind{\cdot=b}=b_R+b_X+b_Y$.
\item \emph{(graft)} Let $R=q_1\cdot q_2$ be split at a vertex $p$, and let $X_1,X_2$ be paths starting
at $p$ with $V(X_1)\cap V(q_1)=\{p\}$ and $V(X_2)\cap V(q_2)=\{p\}$. Then $N_1=q_1\cdot X_1$ and
$N_2=X_2^{-1}\cdot q_2$ are paths with $E(N_1)+E(N_2)=E(R)+E(X_1)+E(X_2)$ and
$b_{N_1}+b_{N_2}+2\ind{\cdot=p}=b_R+b_{X_1}+b_{X_2}$.
\end{enumerate}
\leanst{Gallai.exists\_subst\_list, Gallai.exists\_cut\_edge, Gallai.exists\_graft,
Gallai.exists\_graft\_x}
\end{lemma}

\begin{proof}
(a) Write $R=R_1\cdot ab\cdot R_2$, with $R_1$ ending at $a$ and $R_2$ starting at $b$ (either may
be trivial). As $R$ is a path, $b\notin V(R_1)$ and $a\notin V(R_2)$, so $V(T)\cap V(R_1)=\{a\}$ and
$V(T)\cap V(R_2)=\{b\}$, while $V(R_1)\cap V(R_2)=\emptyset$. Hence $R'=R_1\cdot T\cdot R_2$ is a path;
its ends are the first vertex of $R_1$ and the last vertex of $R_2$, the ends of $R$.
(b) $b\notin V(R_a)$ since $R$ is a path, so $V(R_a)\cap V(X)\subseteq\{a\}$, and $N_1$ is a path
(Lemma~\ref{lem:split}(c)); similarly $a\notin V(R_b)$ gives $V(R_b)\cap V(Y)\subseteq\{b\}$ and $N_2$ is
a path. The ends of $N_1$ are the first vertex of $R_a$ and the far end of $X$; those of $N_2$ are the
far end of $Y$ and the last vertex of $R_b$. The ends of $R,X,Y$ are these four together with $a$ (an
end of $X$) and $b$ (an end of $Y$), which gives the end identity; this includes the case $R_a=a$ or
$R_b=b$, where $a$ or $b$ is also an end of $R$ and remains an end of $N_1$ or $N_2$. The edges of
$N_1,N_2$ are those of $R$ except $ab$, together with those of $X$ and $Y$.
(c) $N_1$ is a path because $V(q_1)\cap V(X_1)=\{p\}$, and $N_2$ because $V(X_2)\cap V(q_2)=\{p\}$.
The ends of $N_1$ are the first vertex of $q_1$ and the far end of $X_1$; those of $N_2$ are the far
end of $X_2$ and the last vertex of $q_2$. The ends of $R,X_1,X_2$ are these four together with the
two ends of $X_1$ and $X_2$ at $p$; if $q_1$ (or $q_2$) is trivial, $p$ is an end of $R$ and of $N_1$
(or $N_2$) as well, and the identity still holds. The edges are those of $R$, $X_1$ and $X_2$.
\end{proof}

When $R$ is a path of a family $L$ and the edge $ab$ (or the vertex $p$) lies on $R$, replacing $R$
in $L$ by $R'$, or by $N_1,N_2$, changes the edge multiset and the end counts as stated.

\subsection{The odd case}

\begin{proposition}[$|C|$ odd]\label{prop:odd}
If $|C|$ is odd, $\RT(B,r)$ holds.
\leanst{Gallai.rootedAt\_oddComp}
\end{proposition}

\begin{proof}
By Lemma~\ref{lem:sides}(c) exactly one port is odd in $A$; by symmetry it is $x$, and $y$ is even
in $A$. Let $|C|=2s+1$ with $s\ge1$. On $C\cup\{x,y\}$ ($2s+3$ vertices) the only even vertex of $A$
is $y$. By Corollary~\ref{cor:lovasz}, $A$ has a decomposition $D_A$ into $s+1$ paths with one end at
each vertex of $C\cup\{x\}$ and no other ends. Let $P\in D_A$ be the path ending at $x$, read from
$x$, and let $c$ be its other end. Then $c\ne x$ has an end, so $c\in C$. The other $s$ paths of
$D_A$ have all their ends in $C$; in particular they have no end at $r$.

\emph{Case 1: $y\notin V(P)$.} Let $W=E+xc$ and $U=(V\setminus C)\cup\{c\}$. The edge $xc$ is not an
edge of $E$ ($c\in C$), $c$ is a leaf of $W$, all edges of $W$ lie in $U$, $W[U]$ is connected by
Lemma~\ref{lem:sides}(a), and $|U|=n-2s$. \emph{Lift:} given $D_W$, let $R$ be its path through
$xc$. Then $V(R)\subseteq U$ and $V(P)\subseteq C\cup\{x\}$, so $V(R)\cap V(P)\subseteq\{x,c\}$.
Substitute $P$ for $xc$ in $R$ (Lemma~\ref{lem:surgery}(a)) and add the other $s$ paths of $D_A$.
The edges are $E(E)+E(A)=E(B)$, the size is $|D_W|+s$, and the ends at $r$ are unchanged, since $R'$
has the ends of $R$ and the added paths have no end at $r$. \emph{Root:} $d_W(r)=d_E(r)+\ind{r=x}$
(as $r\ne c$), while $d_B(r)=d_E(r)+d_A(r)$ with $d_A(r)$ odd for $r=x$, even for $r=y$, and $0$
otherwise. So $d_W(r)\equiv d_B(r)$. Moreover $d_W(r)>0$: for $r=x$ it is at least one, for $r=y$ it
is $d_E(y)>0$, and otherwise $d_W(r)=d_B(r)\ge2$. Lemma~\ref{lem:root} applies.

\emph{Case 2: $y\in V(P)$.} Write $P=T\cdot T'$ with $T$ from $x$ to $y$ and $T'$ from $y$ to $c$;
both are nontrivial, and $V(T)\cap V(T')=\{y\}$. Let $W=E+xy+yc$ on the same $U$; neither edge is in
$E$, so $W$ is a simple graph with $|U|=n-2s$, and $W[U]$ is connected. \emph{Lift:} given $D_W$,
let $R$ be its path through $xy$. Now $V(T)\subseteq(C\cup\{x,y\})\setminus\{c\}$, so
$V(T)\cap V(R)\subseteq V(T)\cap U\subseteq\{x,y\}$; substitute $T$ for $xy$. Let $R_1$ be the path
of the new family through $yc$. Its vertices lie in $U\cup V(T)$, while $V(T')\subseteq(C\cup\{y\})$
avoids $x$; so $V(T')\cap V(R_1)\subseteq\{y,c\}$, using $V(T')\cap V(T)=\{y\}$. Substitute $T'$ for
$yc$ and add the other $s$ paths of $D_A$. As before the edges are
$E(E)+E(T)+E(T')+E(D_A\setminus\{P\})=E(E)+E(A)=E(B)$, the size is $|D_W|+s$, and the ends at $r$ are
unchanged. \emph{Root:}
$d_W(r)=d_E(r)+\ind{r=x}+2\ind{r=y}$, which has the parity of $d_B(r)$ as in Case~1, and is
positive (for $r=y$ it is at least two). Lemma~\ref{lem:root} applies.
\end{proof}

\subsection{The even case with odd ports}

\begin{proposition}[$|C|$ even, both ports odd in $A$]\label{prop:evenodd}
If $|C|$ is even and $d_A(x),d_A(y)$ are odd, $\RT(B,r)$ holds.
\leanst{Gallai.rootedAt\_evenComp\_oddPorts}
\end{proposition}

\begin{proof}
Let $|C|=2s$, $s\ge1$. Every vertex of $C\cup\{x,y\}$ is odd in $A$, so Corollary~\ref{cor:lovasz}
gives $D_A$ with $s+1$ paths and one end at every vertex of $C\cup\{x,y\}$. Let $W=E+xy$ on
$U=V\setminus C$ (the edge $xy$ is not in $E$ by definition). Then $|U|=n-2s$ and $W[U]$ is connected.
\emph{Root:} $d_W(r)=d_E(r)+\ind{r\in\{x,y\}}$ and $d_A(x),d_A(y)$ are odd, so the parities agree;
positivity is as before. \emph{Lift:} let $D_W$ be given and let $R$ be its path through $xy$. Every
path of $D_W$ lies in $U$, and every path of $D_A$ in $C\cup\{x,y\}$, so they meet only in
$\{x,y\}$. Let $X\in D_A$ be the path ending at $x$, read from $x$.

\emph{Case 1: $X$ ends at $y$.} Substitute $X$ for $xy$ in $R$ and add the other $s$ paths of $D_A$,
whose ends lie in $C$. The edges, size and ends at $r$ are as required.

\emph{Case 2: $X$ does not end at $y$.} Let $Y\in D_A$ be the path ending at $y$, read from $y$. Then
$Y\ne X$, $Y$ does not end at $x$ (the only end at $x$ is on $X$), and the other ends of $X$ and $Y$
lie in $C$. Write $R=R_x\cdot xy\cdot R_y$ and cut and attach (Lemma~\ref{lem:surgery}(b)): replace
$R$ by $N_1=R_x\cdot X$ and $N_2=Y^{-1}\cdot R_y$, and add the remaining $s-1$ paths of $D_A$. The
edges are $E(E)+E(A)$ and the size is $|D_W|+1+(s-1)$. At $r\notin C$ we have
$b_X(r)=\ind{r=x}$ and $b_Y(r)=\ind{r=y}$, so
$b_{N_1}(r)+b_{N_2}(r)=b_R(r)$, and the added paths have no end at $r$. Lemma~\ref{lem:root}
applies.
\end{proof}

\subsection{The even case with even ports}
Let $|C|=2s$ with $s\ge1$, and suppose $d_A(x),d_A(y)$ are even. Fix
$w_0\in V\setminus(C\cup\{x,y\})$.

\begin{lemma}[local decompositions]\label{lem:local}
\begin{enumerate}[label=\textup{(\alph*)}]
\item $A$ has a decomposition $D^x$ with at most $s+1$ paths such that $b_{D^x}(v)=2\ind{v=x}$ for
every $v\notin C$.
\item The graph $A'=A+xw_0+w_0y$ has a decomposition $D'$ into $s+1$ paths, all inside
$C\cup\{x,y,w_0\}$, with $b_{D'}(v)=\ind{v=x}+\ind{v=y}$ for every $v\notin C$.
\end{enumerate}
\leanst{Gallai.exists\_localX, Gallai.exists\_localA'}
\end{lemma}

\begin{proof}
(a) Consider $A$ on $C\cup\{x,y\}$ and attach a private leaf at $x$. In $A[C\cup\{x,y\}]^x$
($2s+3$ vertices) the leaf has degree $1$, $x$ has odd degree $d_A(x)+1$, the vertices of $C$ are
odd, and $y$ is the only even vertex. By Theorem~\ref{thm:lovasz} its path number is at most $s+1$.
Lemma~\ref{lem:leaf}(c) gives a decomposition $D_0$ of $A$ with $\lambda_{D_0}(x)\le s+1$, and
Lemma~\ref{lem:ends}(a) ($d_A(x)$ is even and positive) gives $D^x$ with $|D^x|\le s+1$ and
$b(x)\ge2$. Now $2|D^x|=\sum_v b(v)\ge b(x)+b(y)+\sum_{c\in C}b(c)\ge 2+b(y)+2s$, since every
vertex of $C$ is odd and has an end. As $2|D^x|\le2s+2$, we get $b(x)=2$ and $b(y)=0$; vertices
outside $C\cup\{x,y\}$ have degree $0$ in $A$ and no ends.

(b) In $A$ the vertex $w_0$ is isolated, so $xw_0,w_0y$ are new edges. On $C\cup\{x,y,w_0\}$
($2s+3$ vertices) the degrees in $A'$ are $d_A(x)+1$ and $d_A(y)+1$ (odd), $2$ at $w_0$, and odd on
$C$. So $w_0$ is the only even vertex and Corollary~\ref{cor:lovasz} applies.
\end{proof}

\begin{lemma}[the lift]\label{lem:lift-ee}
For every decomposition $D_E$ of $E$ there is a decomposition $D$ of $B$ with $|D|\le|D_E|+s$ and
$b_D(v)=b_{D_E}(v)$ for every $v\notin C$.
\leanst{Gallai.case1\_lift, Gallai.case2\_lift for the family $M$ below, with the ends preserved at
the root $r$; the assembly of $D$ is the step \texttt{finish} in the proof of
Gallai.rootedAt\_evenComp\_evenPorts}
\end{lemma}

\begin{proof}
Since $d_E(x)>0$, some path $R_0\in D_E$ contains $x$; all paths of $D_E$ avoid $C$. It suffices to
find a family $M$ of at most $s+1$ paths with $E(M)=E(R_0)+E(A)$ and $b_M(v)=b_{R_0}(v)$ for every
$v\notin C$: then $D=(D_E\setminus\{R_0\})\cup M$.

\emph{Case 1: $y\notin V(R_0)$.} Take $D^x$ from Lemma~\ref{lem:local}(a). Two distinct paths
$X_1,X_2\in D^x$ end at $x$ (a path has at most one end there); read them from $x$. Their other ends
lie in $C$. Split $R_0=q_1\cdot q_2$ at $x$. Since $V(X_i)\subseteq C\cup\{x,y\}$ and $R_0$ avoids
$C\cup\{y\}$, we have $V(X_i)\cap V(R_0)=\{x\}$. Graft (Lemma~\ref{lem:surgery}(c)) and let $M$
consist of $N_1=q_1\cdot X_1$, $N_2=X_2^{-1}\cdot q_2$ and the at most $s-1$ other paths of $D^x$.
Then $E(M)=E(R_0)+E(A)$, and for $v\notin C$,
$b_M(v)=b_{R_0}(v)+b_{X_1}(v)+b_{X_2}(v)-2\ind{v=x}=b_{R_0}(v)$, because the other paths of $D^x$
have no end outside $C$.

\emph{Case 2: $y\in V(R_0)$.} Take $D'$ from Lemma~\ref{lem:local}(b). The vertex $w_0$ has degree
two in $A'$ and no end, so the path $R\in D'$ through $xw_0$ continues along $w_0y$. Read $R$ so that
\[
R=q\cdot(x\,w_0\,y)\cdot q'',
\]
where $q$ runs from an end of $R$ to $x$ and $q''$ from $y$ to the other end. As $R$ is a path, $q$
avoids $w_0,y$ and $q''$ avoids $w_0,x$, so $V(q)\subseteq C\cup\{x\}$ and $V(q'')\subseteq C\cup\{y\}$,
and $q,q''$ are disjoint. Let $\mathcal R=D'\setminus\{R\}$ ($s$ paths). Both edges at $w_0$ lie on
$R$, so the paths of $\mathcal R$ lie in $C\cup\{x,y\}$, and
\begin{equation}\label{eq:Aedges}
E(\mathcal R)+E(q)+E(q'')=E(A),\qquad b_R(v)+b_{\mathcal R}(v)=\ind{v=x}+\ind{v=y}\ (v\notin C).
\end{equation}
Read $R_0$ so that $x$ comes before $y$: $R_0=q_1\cdot m\cdot q_3$, with $q_1$ from an end of $R_0$ to
$x$, $m$ from $x$ to $y$, and $q_3$ from $y$ to the other end. Then $q_1$ avoids $y$, $q_3$ avoids
$x$, $q_1$ and $q_3$ are disjoint, and all three avoid $C$. Write $t$ for the end of $R$ on $q''$
(so $t=y$ if $q''$ is trivial, and $t\in C$ otherwise).

\emph{Phase $x$.} We produce a path $P=u\cdot q_3$, where $u$ runs from an end of $P$ to $y$, and a
family $M_1$ of $s$ paths, such that:
\begin{enumerate}[label=(P\arabic*),leftmargin=3em]
\item $V(u)\cap V(q'')\subseteq\{y\}$;
\item every path of $M_1$ lies in $V(q_1)\cup C\cup\{x,y\}$;
\item $E(P)+E(M_1)+E(q'')=E(R_0)+E(A)$;
\item $b_P(v)+b_{M_1}(v)+\ind{v=t}=b_{R_0}(v)+\ind{v=y}$ for every $v\notin C$;
\item $b_P(y)\le b_{R_0}(y)$.
\end{enumerate}
\emph{(x-a) $q$ is trivial}, i.e.\ $R$ starts at $x$. Take $P=R_0$ (so $u=q_1\cdot m$) and
$M_1=\mathcal R$. (P1) holds as $u$ avoids $C$; (P2) is clear; (P3) is \eqref{eq:Aedges} with
$E(q)=\emptyset$; (P4) follows from \eqref{eq:Aedges} since $b_R(v)=\ind{v=x}+\ind{v=t}$; (P5) is
trivial.

\emph{(x-b) $q$ is nontrivial.} Its first vertex, an end of $R$, lies in $C$, so $R$ has no end at
$x$ (its other end $t$ is not $x$). By \eqref{eq:Aedges}, $b_{\mathcal R}(x)=1$; let $X\in\mathcal R$
be the path ending at $x$, read from $x$. Graft at $x$ into $R_0=q_1\cdot(m\cdot q_3)$ with $X_1=X$
and $X_2=q^{-1}$ (from $x$ back to the start of $q$). The hypotheses hold: $V(X)\cap V(q_1)=\{x\}$
because $X\subseteq C\cup\{x,y\}$ while $q_1$ avoids $C$ and $y$; and $V(q)\cap V(m\cdot
q_3)=\{x\}$ because $q\subseteq C\cup\{x\}$ while $m\cdot q_3$ avoids $C$. We obtain $N_1=q_1\cdot X$
and $N_2=q\cdot m\cdot q_3$. Put $P=N_2$ with $u=q\cdot m$, and $M_1=\{N_1\}\cup(\mathcal
R\setminus\{X\})$. (P1): $q$ and $q''$ are disjoint, and $m$ meets $q''\subseteq C\cup\{y\}$ only
in $y$. (P2): $N_1\subseteq V(q_1)\cup V(X)$. (P3): $E(N_1)+E(N_2)=E(R_0)+E(X)+E(q)$, and with
\eqref{eq:Aedges} the total is $E(R_0)+E(A)$. (P5): the ends of $P$ are the start of $q$ (in $C$)
and the end of $q_3$, which is an end of $R_0$. (P4): for $v\notin C$, the graft identity gives
$b_{N_1}(v)+b_{N_2}(v)=b_{R_0}(v)+b_X(v)+\ind{v=x}-2\ind{v=x}$, and
$b_{\mathcal R}(v)=\ind{v=x}+\ind{v=y}-\ind{v=t}$ by \eqref{eq:Aedges}, since $R$ starts in $C$. Adding,
$b_P(v)+b_{M_1}(v)=b_{R_0}(v)-\ind{v=x}+b_{\mathcal R}(v)=b_{R_0}(v)+\ind{v=y}-\ind{v=t}$.

\emph{Phase $y$.}
\emph{(y-a) $q''$ is trivial}, i.e.\ $t=y$. Take $M=\{P\}\cup M_1$: it has $s+1$ paths, (P3) gives
$E(M)=E(R_0)+E(A)$, and (P4) gives $b_M(v)=b_{R_0}(v)$ for $v\notin C$.

\emph{(y-b) $q''$ is nontrivial}, so $t\in C$. By (P4) at $v=y$ and (P5), $b_{M_1}(y)\ge1$; let
$Y\in M_1$ end at $y$, read from $y$. Graft at $y$ into $P=u\cdot q_3$ with $X_1=q''$ (attached after
$u$) and $X_2=Y$ (attached before $q_3$). The hypotheses hold: $V(q'')\cap V(u)=\{y\}$ by (P1); and
$V(Y)\cap V(q_3)=\{y\}$ by (P2), since $q_3$ is disjoint from $q_1$ and avoids $C$ and $x$. We obtain
$N_1=u\cdot q''$ and $N_2=Y^{-1}\cdot q_3$; let $M=\{N_1,N_2\}\cup(M_1\setminus\{Y\})$, with $s+1$
paths. Its edges are $E(P)+E(q'')+E(M_1)=E(R_0)+E(A)$ by (P3). For $v\notin C$, using
$b_{q''}(v)=\ind{v=y}$ (as $t\in C$) and (P4) with $\ind{v=t}=0$,
\[
b_M(v)=b_P(v)+\ind{v=y}+b_Y(v)-2\ind{v=y}+b_{M_1}(v)-b_Y(v)=b_{R_0}(v).
\]
This completes the construction of $M$ in all cases.
\end{proof}

The four sub-cases of MR (whether $q$ and $q''$ are trivial, and, when both are nontrivial, whether
the $x$- and $y$-ended paths coincide) are covered by the two optional grafts: the path $Y$ of phase
$y$ may well be the path $N_1$ created in phase $x$, and no separate argument is needed.

\begin{proposition}[$|C|$ even, both ports even in $A$]\label{prop:eveneven}
If $|C|$ is even and $d_A(x),d_A(y)$ are even, $\RT(B,r)$ holds.
\leanst{Gallai.rootedAt\_evenComp\_evenPorts}
\end{proposition}

\begin{proof}
Apply Lemma~\ref{lem:root} with $W=E$ and $U=V\setminus C$, of order $n-2s$; $E[U]$ is connected by
Lemma~\ref{lem:sides}(a). For $r\notin C$, $d_A(r)$ is even (it is $d_A(x)$, $d_A(y)$ or $0$), so
$d_E(r)\equiv d_B(r)$; and $d_E(r)>0$ (Lemma~\ref{lem:sides}(b) at the ports, $d_E(r)=d_B(r)\ge2$
elsewhere). Lemma~\ref{lem:lift-ee} provides the lift.
\end{proof}

\begin{proof}[Proof of Theorem~\ref{thm:sep}]
By Lemma~\ref{lem:sides}(c), either $|C|$ is odd (Proposition~\ref{prop:odd}, after exchanging $x$
and $y$ if necessary), or $|C|$ is even and the ports are both odd (Proposition~\ref{prop:evenodd})
or both even (Proposition~\ref{prop:eveneven}) in $A$.
\end{proof}

\section{The deletion tests on 3-connected graphs: MR Theorem 114}\label{sec:thm114}

In this section we prove Theorem~\ref{thm:main}. Throughout, (T1) and (T2) are as in
Definition~\ref{def:tests} \lean{Gallai.Test1, Gallai.Test2}; assertion (ii) of
Theorem~\ref{thm:main} is \texttt{Gallai.SpineTests} and assertion (iii) is
\texttt{Gallai.SpineTestsNMC}.

\subsection{Gallai implies the tests}

\begin{lemma}[the bridge double]\label{lem:bridge}
Assume Gallai's conjecture. If $H$ is connected of even order $N$ and $r\in V(H)$, then
$p(H^r)\le N/2$.
\leanst{Gallai.dbl, Gallai.two\_mul\_pathNumber\_addLeaf\_le, Gallai.rootedTest\_of\_gallai}
\end{lemma}

\begin{proof}
Let $\Delta$ consist of two disjoint copies $H_1,H_2$ of $H$, with copies $r_1,r_2$ of $r$, joined
by the edge $r_1r_2$. It is connected of order $2N$, so $p(\Delta)\le N$. Let $D$ be a decomposition
of $\Delta$ of size $p(\Delta)$. The bridge $r_1r_2$ lies on exactly one path $Q$, and every other
path lies in $H_1$ or in $H_2$; let $D_i$ be the paths in $H_i$, so $|D_1|+|D_2|+1=|D|$. Write
$Q=Q_1\cdot r_1r_2\cdot Q_2$ with $Q_i$ in $H_i$ (possibly trivial). Then $D_1\cup\{Q_1\cdot
r_1r_2\}$ is a decomposition of $H_1$ with the private leaf $r_2$ at $r_1$, a copy of $H^r$, so
$p(H^r)\le|D_1|+1$; likewise $p(H^r)\le|D_2|+1$. Hence $2p(H^r)\le|D|+1\le N+1$, and
$p(H^r)\le\fl{(N+1)/2}=N/2$.
\end{proof}

\begin{lemma}[rooted bounds from Gallai]\label{lem:gallai-rooted}
Assume Gallai's conjecture. Let $H$ be connected of order $N$ and $r\in V(H)$.
\begin{enumerate}[label=\textup{(\alph*)}]
\item If $N$ is even and $d_H(r)$ is even and positive, some decomposition $D$ of $H$ has
$|D|\le N/2$ and $b_D(r)\ge2$.
\item If $N$ is odd and $d_H(r)$ is odd and at least $3$, some $D$ has $|D|\le(N+1)/2$ and
$b_D(r)\ge3$.
\end{enumerate}
\leanst{the arguments of Gallai.test1\_of\_gallai and Gallai.test2\_of\_gallai, stated there for
$H=J-z$}
\end{lemma}

\begin{proof}
(a) By Lemma~\ref{lem:bridge}, $p(H^r)\le N/2$; Lemma~\ref{lem:leaf}(c) gives $D$ with
$\lambda_D(r)\le N/2$, and Lemma~\ref{lem:ends}(a) finishes. (b) The graph $H^r$ is connected of even
order $N+1$, and $r$ has even positive degree $d_H(r)+1$ in it. By (a), $H^r$ has a decomposition
with at most $(N+1)/2$ paths and two ends at $r$. Trimming the leaf (Lemma~\ref{lem:leaf}(a)) gives a
decomposition $D$ of $H$ with either $|D|\le(N+1)/2$ and $b_D(r)\ge3$, or $|D|\le(N-1)/2$ and
$b_D(r)\ge1$. In the second case Lemma~\ref{lem:ends}(b) with $k=(N-1)/2$ gives at most $(N+1)/2$
paths and three ends at $r$.
\end{proof}

\begin{proposition}[Gallai implies the tests]\label{prop:forward}
If Gallai's conjecture holds, then \textup{(T1)} and \textup{(T2)} hold for every $3$-connected $J$
and every $z\in V(J)$, whether or not $J-z$ has a matching cut.
\leanst{Gallai.spineTests\_of\_gallai, Gallai.test1\_of\_gallai, Gallai.test2\_of\_gallai}
\end{proposition}

\begin{proof}
Let $n=|V(J)|$. By Lemma~\ref{lem:kconn}, $J-z$ is connected and $d_J(v)\ge3$ for every $v$. For a
neighbour $r$ of $z$ we have $d_{J-z}(r)=d_J(r)-1$ \lean{Gallai.delVert\_degree}.

(T1) If $z$ is the only even vertex, the $n-1$ odd vertices are even in number, so $n$ is odd and
$J-z$ has even order $n-1$. The vertex $r$ is odd in $J$ with $d_J(r)\ge3$, so $d_{J-z}(r)$ is even
and at least two. Lemma~\ref{lem:gallai-rooted}(a) gives at most $(n-1)/2$ paths with two ends at
$r$.

(T2) If $z,r$ are the only even vertices, then $n-2$ is even, so $J-z$ has odd order $n-1$. Now
$d_J(r)$ is even and at least three, hence at least four, so $d_{J-z}(r)$ is odd and at least three.
Lemma~\ref{lem:gallai-rooted}(b) gives at most $n/2$ paths with three ends at $r$.
\end{proof}

By Remark~\ref{rem:59conv}, Lemma~\ref{lem:gallai-rooted} also shows that Gallai's conjecture
implies $\RT(B,x)$ for all $2$-connected $B$.

\subsection{The parity apex}

\begin{lemma}[apex lemma; A2 \S4.2]\label{lem:apex}
Let $B$ be $2$-connected and $S\subseteq V(B)$ with $|S|\ge3$. Suppose that for all distinct
$x,y\in V(B)$, every vertex $v\notin\{x,y\}$ is joined to a vertex of $S\setminus\{x,y\}$ by a path of
$B-\{x,y\}$ (equivalently, every component of $B-\{x,y\}$ meets $S$). Then $J=B+z(S)$ is
$3$-connected.
\leanst{Gallai.apex\_kConnected\_three}
\end{lemma}

\begin{proof}
$|V(J)|=|V(B)|+1\ge4$. Let $X\subseteq V(J)$ with $|X|\le2$; we show $J-X$ is connected. If $z\in X$,
then $J-X=B-(X\setminus\{z\})$ with $|X\setminus\{z\}|\le1$, which is connected because $B$ is
$2$-connected. If $z\notin X$, it suffices to join every $v\in V(B)\setminus X$ to $z$ in $J-X$. If
$|X|\le1$, pick $t\in S\setminus X$ (possible as $|S|\ge3$); $B-X$ is connected and contains a path
from $v$ to $t$, which we extend by the edge $tz$. If $X=\{x,y\}$, the hypothesis gives a path of
$B-\{x,y\}$ from $v$ to some $s\in S\setminus\{x,y\}$, extended by $sz$.
\end{proof}

\subsection{The tests imply Gallai}

\begin{lemma}[few even vertices]\label{lem:fewEven}
Let $B$ be $2$-connected of order $n$, let $r\in V(B)$, and suppose $S=\Ev(B)\cup\{r\}$ has at most
two elements. Then $\RT(B,r)$ holds.
\leanst{Gallai.rootedAt\_of\_few\_even}
\end{lemma}

\begin{proof}
$d_B(r)\ge2$ by Lemma~\ref{lem:kconn}. (R1): let $n$ and $d_B(r)$ be even. In $B^r$ the leaf has
degree $1$ and $r$ has odd degree $d_B(r)+1$; the even vertices of $B^r$ are those of
$\Ev(B)\setminus\{r\}=S\setminus\{r\}$, at most one. By Theorem~\ref{thm:lovasz},
$p(B^r)\le\fl{(n+1)/2}=n/2$. Lemma~\ref{lem:leaf}(c) and Lemma~\ref{lem:ends}(a) give at most $n/2$
paths with two ends at $r$. (R2): let $n$ and $d_B(r)$ be odd. Then $\Ev(B)\subseteq
S\setminus\{r\}$ has at most one element, so $p(B)\le\fl{n/2}$ by Theorem~\ref{thm:lovasz}. As
$d_B(r)$ is odd and at least two, it is at least three, and Lemma~\ref{lem:ends}(b) gives at most
$\fl{n/2}+1$ paths with three ends at $r$.
\end{proof}

\begin{proposition}[the deletion tests give the rooted tests; MR p.~129]\label{prop:rootedFromTests}
Assume assertion \textup{(iii)} of Theorem~\ref{thm:main}. Let $B\in\cC$ have order $n$ and assume
$(\MBU_{n-2})$. Then $\RT(B,r)$ holds for every $r\in V(B)$.
\leanst{Gallai.rootedAt\_of\_tests}
\end{proposition}

\begin{proof}
Let $S=\Ev(B)\cup\{r\}$. If $|S|\le2$, use Lemma~\ref{lem:fewEven}. If some separated component $C$
for some pair $\{x,y\}$ has $r\notin C$ and only odd vertices, use Theorem~\ref{thm:sep}. Assume
neither happens; so $|S|\ge3$.

We check the hypothesis of Lemma~\ref{lem:apex} \lean{Gallai.exists\_sepComponent\_or}. Let
$x\ne y$, $v\notin\{x,y\}$, and let $K$ be the vertex set of the component of $B-\{x,y\}$ containing
$v$. If $K$ meets $S$, it contains a vertex of $S\setminus\{x,y\}$ joined to $v$ inside $K$.
Otherwise $K$ is nonempty and connected, every neighbour of a vertex of $K$ lies in $K\cup\{x,y\}$,
and some vertex of $S\setminus\{x,y\}$ (which exists as $|S|\ge3$) lies outside $K\cup\{x,y\}$; so
$K$ is a separated component for $\{x,y\}$. It avoids $r\in S$, and all its vertices are odd since
$\Ev(B)\subseteq S$. This is the excluded case. Hence $J=B+z(S)$ is $3$-connected by
Lemma~\ref{lem:apex}. Moreover $J-z=B$ has no matching cut, as $B\in\cC$
\lean{Gallai.not\_hasMatchingCut\_delVert\_apex}.

The degrees in $J$ are $d_J(z)=|S|$, $d_J(v)=d_B(v)+1$ for $v\in S$ and $d_J(v)=d_B(v)$ otherwise.
Let $o$ be the number of odd vertices of $B$, which is even.

(R1) Let $n$ and $d_B(r)$ be even. Then $r\in\Ev(B)$, $S=\Ev(B)$ and $|S|=n-o$ is even. So $z$ is
even, every vertex of $S$ becomes odd, and every other vertex stays odd: $z$ is the only even vertex
of $J$, and $r$ is a neighbour of $z$. By (T1), $J-z=B$ has a decomposition into at most
$(|V(J)|-1)/2=n/2$ paths with two ends at $r$.

(R2) Let $n$ and $d_B(r)$ be odd. Then $r\notin\Ev(B)$ and $|S|=(n-o)+1$ is even, since $n-o$ is odd.
So $z$ is even, $r$ becomes even, the vertices of $\Ev(B)$ become odd, and all others stay odd: $z$
and $r$ are the only even vertices of $J$, and they are adjacent. By (T2), $B$ has a decomposition
into at most $|V(J)|/2=(n+1)/2=\fl{n/2}+1$ paths with three ends at $r$.
\end{proof}

Since $|S|$ is even in both cases, $|S|\ge4$, so $d_J(z)\ge4$.

\begin{theorem}[the tests imply Gallai]\label{thm:converse}
Assertion \textup{(iii)} of Theorem~\ref{thm:main} implies Gallai's conjecture.
\leanst{Gallai.gallai\_of\_spineTestsNMC}
\end{theorem}

\begin{proof}
We show by strong induction on $n$ that every $B\in\cC$ of order $n$ satisfies $\RT(B,r)$ for every
$r$. Let $B\in\cC$ have order $n$ (so $n\ge3$), and assume the claim for all orders below $n$; this
is $(\RT_{n-1})$. Theorem~\ref{thm:59} with $m=n-2$ gives $(\MBU_{n-2})$, and
Proposition~\ref{prop:rootedFromTests} gives $\RT(B,r)$.

Hence $(\RT_N)$ holds for every $N$. Let $G$ be connected of order $N$ and $r\in V(G)$.
Theorem~\ref{thm:59} with $m=N$ gives $(\MBU_N)$, which applies to $G$ because its $2$-connected
blocks have order at most $N$. So $q(G,\{r\})\le\fl{(N-1)/2}$, and by
Lemma~\ref{lem:marked-basic}(a), $p(G)\le\fl{(N-1)/2}+1=\cl{N/2}$.
\end{proof}

\begin{proof}[Proof of Theorem~\ref{thm:main}]
(i)$\Rightarrow$(ii) is Proposition~\ref{prop:forward}, (ii)$\Rightarrow$(iii) is trivial, and
(iii)$\Rightarrow$(i) is Theorem~\ref{thm:converse}. The ingredients are Lov\'asz's theorem
(Theorem~\ref{thm:lovasz}), Theorem~59 (Theorem~\ref{thm:59}) and Theorem~113
(Theorem~\ref{thm:sep}) \lean{Gallai.gallai\_iff\_spineTests, Gallai.gallai\_iff\_spineTests'}.
\end{proof}

\section{Cones and rooted Gallai for Eulerian graphs}\label{sec:cones}

For a graph $B$ with vertex set $W$, the \emph{cone} $J=B+z$ is the graph $B+z(W)$ of
Section~\ref{sec:prelim}: the new vertex $z$ (the \emph{apex}) is joined to every vertex of $W$, and
$J-z=B$ \lean{Gallai.fullCone}. Write $n=|W|$. For $r\in W$ we write $\mathrm{RD}(B,r)$ for the
assertion that $B$ has a decomposition $D$ with $|D|\le n/2$ and $b_D(r)\ge2$
\lean{Gallai.RootedDecomp}. When $d_B(r)$ is even and positive, $\mathrm{RD}(B,r)$ is equivalent
to $\pr(B)\le n/2$, because $b_D(r)\ge1$ forces $b_D(r)\ge2$ by Lemma~\ref{lem:parity}. So
$\RE(B,r)$ of Definition~\ref{def:RE} is $\mathrm{RD}(B,r)$ for $2$-connected Eulerian $B$ of even
order. In the formalisation the apex is a new element of the vertex type, and $J-z$ is isomorphic
rather than equal to $B$ \lean{Gallai.fullConeDelIso}; decompositions are transported along the
isomorphism with their sizes and end counts \lean{Gallai.PathDecomp.exists\_of\_iso}.

\begin{proposition}\label{prop:cone-par}
$d_J(z)=n$ and $d_J(a)=d_B(a)+1$ for $a\in W$. Hence $z$ is the only even vertex of $J$ if and only
if $n$ is even and $B$ is Eulerian.
\leanst{Gallai.fullCone\_degree\_none, Gallai.fullCone\_degree\_some, Gallai.fullCone\_onlyEven\_iff}
\end{proposition}

\begin{proof}
$N_J(z)=W$ and $N_J(a)=N_B(a)\cup\{z\}$. So $d_J(a)$ is odd if and only if $d_B(a)$ is even.
\end{proof}

\begin{proposition}\label{prop:cone-conn}
$J$ is $3$-connected if and only if $B$ is $2$-connected.
\leanst{Gallai.kConnected\_three\_fullCone\_iff}
\end{proposition}

\begin{proof}
Suppose $B$ is $2$-connected and $X\subseteq V(J)$ has $|X|\le 2$. If $z\in X$, then
$J-X=B-(X\setminus\{z\})$ is connected, since $|X\setminus\{z\}|\le1$. If $z\notin X$, then $z$
survives and is adjacent to every other vertex of $J-X$, so $J-X$ is connected. Also
$|V(J)|=n+1\ge4$. Conversely, suppose $J$ is $3$-connected and $X\subseteq W$ has $|X|\le1$. Then
$J-(X\cup\{z\})=B-X$ is connected. Also $n=|V(J)|-1\ge3$.
\end{proof}

\begin{proposition}\label{prop:cone-mc}
$J-z$ has a matching cut if and only if $B$ has one.
\leanst{Gallai.hasMatchingCut\_delVert\_fullCone\_iff, Gallai.HasMatchingCut.of\_iso}
\end{proposition}

\begin{proof}
$J-z=B$.
\end{proof}

\begin{proposition}\label{prop:cone-t1}
For every $r\in W$, the conclusion of \textup{(T1)} at $(J,z,r)$ holds if and only if
$\mathrm{RD}(B,r)$ holds. Consequently \textup{(T1)} holds at $(J,z)$ if and only if: when $n$ is
even and $B$ is Eulerian, $\mathrm{RD}(B,r)$ holds for every $r\in W$.
\leanst{Gallai.test1At\_fullCone\_iff, Gallai.test1\_fullCone\_iff}
\end{proposition}

\begin{proof}
The conclusion of (T1) asks for a decomposition of $J-z=B$ into at most $(|V(J)|-1)/2=n/2$ paths
with $b(r)\ge2$, which is $\mathrm{RD}(B,r)$. The neighbours of $z$ are the vertices of $W$, and the
hypothesis of (T1) is described by Proposition~\ref{prop:cone-par}.
\end{proof}

No connectivity is used in Proposition~\ref{prop:cone-t1}.

\begin{proposition}\label{prop:cone-vac}
The other tests at a cone are vacuous: \textup{(T2)} at $z$ always holds; \textup{(T1)} at every
$a\in W$ always holds; and \textup{(T2)} at every $a\in W$ holds when $n$ is even.
\leanst{Gallai.test2\_fullCone\_apex, Gallai.test1\_fullCone\_some,
Gallai.test2\_fullCone\_some\_of\_even, Gallai.not\_unique\_odd}
\end{proposition}

\begin{proof}
No graph has exactly one odd vertex, by the handshake lemma. We show that the hypothesis of each test
fails.
\begin{itemize}[leftmargin=*]
\item (T2) at $(z,a)$ needs $z$ and $a$ to be the only even vertices of $J$. By
Proposition~\ref{prop:cone-par}, $a$ would then be the only odd vertex of $B$.
\item (T1) at $a$ needs $a$ to be the only even vertex of $J$. Then $z$ is odd, so $n$ is odd and
$|V(J)|$ is even, and $J$ would have the odd number $|V(J)|-1$ of odd vertices.
\item (T2) at $(a,x)$ with $n$ even: $z$ is even and $z\ne a$, so $x=z$, and $a$ would be the only
odd vertex of $B$.\qedhere
\end{itemize}
\end{proof}

When $n$ is odd, (T2) at a vertex $a\in W$ is a genuine test. It concerns graphs $B$ with exactly two
odd vertices, which are adjacent, and it is not of the form RE.

\begin{proposition}[the leaf form]\label{prop:cone-leaf}
Let $r\in W$ with $d_B(r)$ even and positive. Then $\mathrm{RD}(B,r)$ holds if and only if
$p(B^r)\le n/2$.
\leanst{Gallai.rootedDecomp\_iff\_addLeaf}
\end{proposition}

\begin{proof}
If $D$ witnesses $\mathrm{RD}(B,r)$, then $b_D(r)\ge1$, so $p(B^r)\le\lambda_D(r)=|D|\le n/2$ by
Lemma~\ref{lem:leaf}(b). Conversely, Lemma~\ref{lem:leaf}(c) gives a decomposition $D$ of $B$ with
$\lambda_D(r)\le p(B^r)\le n/2$, and Lemma~\ref{lem:ends}(a) with $k=\fl{n/2}$ gives $\mathrm{RD}(B,r)$.
\end{proof}

\begin{theorem}[Theorem~\ref{thm:cones-intro}]\label{thm:cones}
Let $B$ be a graph of order $n$ and $J=B+z$.
\begin{enumerate}[label=\textup{(\alph*)}]
\item $J$ is $3$-connected if and only if $B$ is $2$-connected, and $J-z$ has a matching cut if and
only if $B$ has one.
\item \textup{(T1)} holds at $(J,z)$ if and only if: when $n$ is even and $B$ is Eulerian,
$\mathrm{RD}(B,r)$ holds for every $r\in W$. \textup{(T2)} holds at $(J,z)$.
\item Assertion \textup{(iii)} of Theorem~\ref{thm:main} implies $\RE(B,r)$ for every $2$-connected
Eulerian graph $B$ of even order without a matching cut and every $r\in V(B)$. Conversely, if $n$ is
even, then \textup{(T1)} and \textup{(T2)} hold at every vertex of $B+z$ if and only if
$\mathrm{RD}(B,r)$ holds for every $r$ whenever $B$ is Eulerian.
\item Gallai's conjecture implies $\mathrm{RD}(B,r)$ for every connected Eulerian graph $B$ of even
order and every vertex $r$ of positive degree.
\end{enumerate}
\leanst{Gallai.rootedDecomp\_of\_spineTestsNMC, Gallai.fullConeTests\_iff, Gallai.rootedDecomp\_of\_gallai}
\end{theorem}

\begin{proof}
(a) is Propositions~\ref{prop:cone-conn} and~\ref{prop:cone-mc}, and (b) is
Propositions~\ref{prop:cone-t1} and~\ref{prop:cone-vac}. (c) Let $B$ be $2$-connected and Eulerian
of even order without a matching cut. By (a), $J$ is $3$-connected and $J-z$ has no matching cut, so
(iii) gives (T1) at $(J,z)$, and (b) gives $\mathrm{RD}(B,r)$ for every $r$. For the converse
statement, let $n$ be even. By Proposition~\ref{prop:cone-vac} the tests at the vertices of $W$
hold, and so does (T2) at $z$; by (b), (T1) at $z$ is the stated condition. (d) Apply
Lemma~\ref{lem:gallai-rooted}(a) with $H=B$; here $d_B(r)$ is even and positive.
\end{proof}

Theorem~\ref{thm:cones} locates RE inside the spine exactly. Among the instances of
Theorem~\ref{thm:main}(iii) at the apex of a cone, the ones whose hypotheses can hold are precisely
the statements $\RE(B,r)$ for $2$-connected Eulerian graphs $B$ of even order without a matching
cut. At the other vertices of a cone over a graph of even order all tests are vacuous. Over a graph
of odd order, (T2) at a non-apex vertex is a genuine test that is not of the form RE.
Theorem~\ref{thm:main} and Theorem~\ref{thm:cones}(c),(d) give
\[
\text{Gallai}\iff\text{(iii)}\implies\RE\text{ on $2$-connected Eulerian graphs of even order
with no matching cut,}
\]
while Gallai's conjecture also gives RE in full, and RE gives Gallai's bound $p(B)\le\pr(B)\le n/2$ for
$2$-connected Eulerian graphs of even order.

\begin{example}\label{ex:K26}
Cycles of length at least $4$ have matching cuts: split the cycle into two paths. The complete
bipartite graph $K_{2,6}$ is $2$-connected and Eulerian of order $8$, and it has no matching cut. Let
$\{a,b\}$ be its side of size two, and let $(X,Y)$ be a partition of the vertex set into nonempty
sets. If $a,b\in X$, then a vertex of $Y$ (which has degree $2$) has two neighbours in $X$. If
$a\in X$ and $b\in Y$, then $a$ has six neighbours, and if at most one of them lies in $Y$, then at
least five lie in $X$, where each is adjacent to $b\in Y$; so $b$ has at least five neighbours in
$X$. In either case some vertex has two neighbours on the other side. So $\RE(K_{2,6},r)$ is one
of the instances that the spine requires.
\end{example}

\section{Using the reduction}\label{sec:use}

\subsection*{The smallest counterexample}
The proof of Theorem~\ref{thm:converse} gives more than the equivalence: it shows what a smallest
counterexample looks like, and which inductive hypotheses are available when one tries to prove a
test.

\begin{corollary}[shape of a smallest counterexample]\label{cor:mincex}
Suppose Gallai's conjecture fails. Then there are a graph $B\in\cC$ of some order $n$ and a vertex
$r$ such that $\RT(B,r)$ fails and $(\RT_{n-1})$ holds. For any such $B$ and $r$, put
$S=\Ev(B)\cup\{r\}$ and $J=B+z(S)$. Then:
\begin{enumerate}[label=\textup{(\alph*)}]
\item the marked bound $(\MBU_{n-2})$ holds, so every connected graph whose $2$-connected blocks have
order at most $|V(J)|-3$ satisfies Gallai's bound;
\item $|S|$ is even and at least $4$, $J$ is $3$-connected, and $J-z=B$ has no matching cut;
\item no separated component for any pair avoids $r$ and consists of odd vertices only;
\item if $n$ is even, then $z$ is the only even vertex of $J$ and \textup{(T1)} fails at $(J,z)$
with root $r$; if $n$ is odd, then $z$ and $r$ are the only even vertices of $J$ and \textup{(T2)}
fails at $(J,z)$ with root $r$.
\end{enumerate}
\end{corollary}

\begin{proof}
If $\RT(B,r)$ held for all $B\in\cC$ and all $r$, the second paragraph of the proof of
Theorem~\ref{thm:converse} would give Gallai's conjecture. So some $B\in\cC$ fails; choose one of
least order $n$, with a root $r$ at which it fails. Then $(\RT_{n-1})$ holds by minimality.

Now let $B$ and $r$ be any such pair. (a) Theorem~\ref{thm:59} with $m=n-2$ gives $(\MBU_{n-2})$,
and Lemma~\ref{lem:marked-basic}(a) with one mark gives Gallai's bound. (c) Otherwise
Theorem~\ref{thm:sep}, which applies by (a), would give $\RT(B,r)$. (b) By Lemma~\ref{lem:fewEven},
$|S|\ge3$. The second paragraph of the proof of Proposition~\ref{prop:rootedFromTests} uses only
(c), $|S|\ge3$ and $B\in\cC$, and shows that $J$ is $3$-connected with $J-z=B$ free of matching cuts. $\RT(B,r)$
fails, so one of (R1), (R2) fails. If (R1) fails, $n$ and $d_B(r)$ are even, and the third paragraph
of that proof shows that $|S|$ is even, that $z$ is the only even vertex of $J$, and that the
conclusion of (T1) at root $r$ is exactly (R1); so (T1) fails. If (R2) fails, $n$ and $d_B(r)$ are
odd, and in the same way $|S|$ is even, $z,r$ are the only even vertices, and the conclusion of
(T2) at $r$ is (R2). Since $|S|\ge3$ is even, $|S|\ge4$. This proves (b) and (d).
\end{proof}

So to prove Gallai's conjecture it suffices to prove (T1) and (T2) at the \emph{parity apexes}
$J=B+z(S)$, $S=\Ev(B)\cup\{r\}$, of graphs $B\in\cC$ without an all-odd separated component avoiding
$r$, and one may assume the marked bound for all graphs with smaller blocks. Corollary~\ref{cor:mincex}
is not stated as a separate Lean declaration; its steps are those of \texttt{rootedAt\_of\_tests}
and \texttt{gallai\_of\_spineTestsNMC}.

\subsection*{Lov\'asz decompositions as a starting point}
In (T1), $J$ has a single even vertex $z$, so Theorem~\ref{thm:lovasz-strong} decomposes $J$
itself into $(|V(J)|-1)/2$ paths with one end at every odd vertex. The test asks for the same number
of paths in $J-z$, with two ends at $r\in N(z)$. In (T2), deleting the edge $zr$ leaves a graph
with no even vertex, which Theorem~\ref{thm:lovasz-strong} decomposes into $|V(J)|/2$ paths with one
end at every vertex. This suggests a programme: delete a set of edges at $z$ (and possibly further
edges among its neighbours), decompose the remainder with a Lov\'asz decomposition whose ends are
constrained, and lift the result back to $J-z$ without increasing the number of paths. Fan's lemma
(Lemma~\ref{lem:fan}) is the basic move for such lifts. The programme is carried out in companion
papers, which we cite by title only and on which nothing here depends:
\begin{itemize}
\item \emph{Constrained Lov\'asz decompositions: separation, neighbour avoidance, leaf tails and
prescribed ends}~\cite{P2}, on strengthenings of Theorem~\ref{thm:lovasz-strong} with constraints on
the ends;
\item \emph{Reversible lifts from constrained Lov\'asz decompositions to the spine tests}~\cite{P3},
on reducing (T1) and (T2) to constrained Lov\'asz statements;
\item \emph{Rooted Gallai for Eulerian graphs, and the two-end test at near-universal roots}~\cite{P4},
on RE and on (T1) at roots of large degree.
\end{itemize}

\subsection*{The cone case}
By Theorem~\ref{thm:cones}, any proof of (T1) must in particular prove RE for $2$-connected
Eulerian graphs of even order without a matching cut, such as $K_{2,6}$. Conversely, by
Theorem~\ref{thm:cones}(d), a single connected Eulerian graph $B$ of even order and a vertex $r$ of
positive degree with $\pr(B)>|V(B)|/2$ would disprove Gallai's conjecture.

\subsection*{Formal use}
In the Lean library, a proof of the proposition \texttt{SpineTestsNMC} (a closed statement about all
finite graphs in a universe) yields, through \texttt{gallai\_iff\_spineTests'}, a proof of
\texttt{GallaiConjecture} in the same universe. Partial results can be stated against the predicates
\texttt{Test1}, \texttt{Test2} and \texttt{Test1At} of Section~\ref{sec:lean}, and on cones against
\texttt{RootedDecomp}.

\section{The formalisation}\label{sec:lean}

\subsection{What is checked}
Every result of Sections~\ref{sec:prelim}--\ref{sec:cones} that is stated as a theorem,
proposition, lemma or corollary is proved in the Lean~4 library \texttt{GallaiLean} of the project
repository (directory \texttt{lean/}), on top of Mathlib, with four exceptions:
\begin{itemize}
\item Lemma~\ref{lem:gallai-rooted}(a) is formalised in general, as the combination of
\texttt{rootedTest\_of\_gallai} and \texttt{rootedDecomp\_iff\_addLeaf}; part (b) only for $H=J-z$,
inside \texttt{test2\_of\_gallai};
\item Lemma~\ref{lem:lift-ee} is formalised with the ends preserved at the chosen root;
\item Remark~\ref{rem:59conv} (Gallai implies the rooted tests) and Corollary~\ref{cor:mincex} are not
formalised as separate statements;
\item Example~\ref{ex:K26} is not formalised.
\end{itemize}
The toolchain is \texttt{leanprover/lean4:v4.36.0} with Mathlib at tag \texttt{v4.36.0} (commit
\texttt{fc4f59e577feb1a4956671bc04ac18aaaf8cd842}, from \texttt{lean/lake-manifest.json}). The
library contains no \texttt{sorry}. Lov\'asz's theorem is not an assumption: it is proved in
\texttt{Lovasz.lean} as in Section~\ref{sec:lovasz}, from Fan's lemma, which is proved in
\texttt{Fan.lean} and \texttt{FanMoves.lean}. Theorem~\ref{thm:main} is the conjunction
{\footnotesize
\begin{verbatim}
theorem gallai_iff_spineTests' :
    (GallaiConjecture.{u} ↔ SpineTests.{u}) ∧ (GallaiConjecture.{u} ↔ SpineTestsNMC.{u}) :=
  gallai_iff_spineTests thm59
\end{verbatim}
}
in \texttt{Separator.lean}, where \texttt{gallai\_iff\_spineTests} combines
\texttt{spineTests\_of\_gallai} with \texttt{gallai\_of\_spineTestsNMC lovaszStatement h59
separatorStatement}. Several cited declarations take Lov\'asz's theorem, Theorem~59 or
Theorem~113 as explicit hypotheses (\texttt{LovaszStatement}, \texttt{Thm59Statement},
\texttt{SeparatorStatement}): for example \texttt{separator\_theorem},
\texttt{rootedAt\_of\_few\_even}, \texttt{rootedAt\_of\_tests} and
\texttt{gallai\_of\_spineTestsNMC}. These hypotheses are discharged by \texttt{lovaszStatement},
\texttt{thm59} and \texttt{separatorStatement}, as in the line above. The statement is
universe-polymorphic: for every universe \texttt{u} it relates the conjecture for all finite vertex
types in \texttt{Type u} to the tests for all finite vertex types in \texttt{Type u}.

\subsection{The definitions}
The definitions behind the statement are the following (\texttt{Basic.lean},
\texttt{ThreeConnected.lean} and \texttt{EulerianCone.lean}; \texttt{V} is a type with
\texttt{[Fintype V] [DecidableEq V]}). Declarations made inside a namespace are shown with their
full names, and the comments are ours; otherwise the text is verbatim.
{\footnotesize
\begin{verbatim}
structure GPath (G : SimpleGraph V) where          -- a path: a simple walk of length ≥ 1
  src : V
  tgt : V
  walk : G.Walk src tgt
  isPath : walk.IsPath
  pos : 0 < walk.length
def GPath.edges (P : GPath G) : List (Sym2 V) := P.walk.edges
def GPath.endsAt [DecidableEq V] (P : GPath G) (v : V) : ℕ :=
  (if P.src = v then 1 else 0) + (if P.tgt = v then 1 else 0)
structure PathDecomp (G : SimpleGraph V) where     -- edge lists partition E(G)
  paths : List (GPath G)
  nodup_edges : (paths.flatMap GPath.edges).Nodup
  covers : ∀ e, e ∈ G.edgeSet ↔ e ∈ paths.flatMap GPath.edges
def PathDecomp.size (D : PathDecomp G) : ℕ := D.paths.length
def PathDecomp.ends [DecidableEq V] (D : PathDecomp G) (v : V) : ℕ :=
  (D.paths.map (fun P => P.endsAt v)).sum
noncomputable def pathNumber (G : SimpleGraph V) : ℕ :=
  sInf {k | ∃ D : PathDecomp G, D.size = k}
def GallaiConjecture : Prop :=
  ∀ (W : Type u) [Fintype W] (G : SimpleGraph W), G.Connected →
    pathNumber G ≤ (Fintype.card W + 1) / 2

def KConnected (k : ℕ) (G : SimpleGraph V) : Prop :=
  k + 1 ≤ Fintype.card V ∧ ∀ X : Finset V, X.card < k → (G.induce ((X : Set V)ᶜ)).Connected
abbrev delVert (J : SimpleGraph V) (z : V) : SimpleGraph {v | v ≠ z} := J.induce {v | v ≠ z}
def HasMatchingCut (G : SimpleGraph V) : Prop :=
  ∃ X : Set V, X.Nonempty ∧ Xᶜ.Nonempty ∧
    ∀ a b c, G.Adj a b → G.Adj a c → (a ∈ X ↔ b ∉ X) → (a ∈ X ↔ c ∉ X) → b = c
def Test1 (J : SimpleGraph V) [DecidableRel J.Adj] (z : V) : Prop :=
  (∀ v, Even (J.degree v) ↔ v = z) → ∀ r (hr : J.Adj z r),
    ∃ D : PathDecomp (delVert J z),
      D.size ≤ (Fintype.card V - 1) / 2 ∧ 2 ≤ D.ends ⟨r, hr.ne'⟩
def Test2 (J : SimpleGraph V) [DecidableRel J.Adj] (z : V) : Prop :=
  ∀ r (hr : J.Adj z r), (∀ v, Even (J.degree v) ↔ v = z ∨ v = r) →
    ∃ D : PathDecomp (delVert J z), D.size ≤ Fintype.card V / 2 ∧ 3 ≤ D.ends ⟨r, hr.ne'⟩
def SpineTests : Prop :=
  ∀ (V : Type u) [Fintype V] [DecidableEq V] (J : SimpleGraph V) [DecidableRel J.Adj] (z : V),
    KConnected 3 J → Test1 J z ∧ Test2 J z
def SpineTestsNMC : Prop :=
  ∀ (V : Type u) [Fintype V] [DecidableEq V] (J : SimpleGraph V) [DecidableRel J.Adj] (z : V),
    KConnected 3 J → ¬ HasMatchingCut (delVert J z) → Test1 J z ∧ Test2 J z

def fullCone (B : SimpleGraph W) : SimpleGraph (Option W)   -- the cone B + z, apex `none`
def RootedDecomp (B : SimpleGraph W) (r : W) : Prop :=
  ∃ D : PathDecomp B, D.size ≤ Fintype.card W / 2 ∧ 2 ≤ D.ends r
\end{verbatim}
}
These match Section~\ref{sec:intro} and Definition~\ref{def:tests} exactly. Natural-number division
is the floor, so $(|V|+1)/2=\lceil|V|/2\rceil$; in \texttt{Test1} the order is odd and in
\texttt{Test2} it is even, so the budgets $(n-1)/2$ and $n/2$ are exact. \texttt{D.ends v} counts
path ends at $v$ with multiplicity, which is $b_D(v)$. In \texttt{HasMatchingCut} the last clause says
that every vertex has at most one neighbour on the other side. \texttt{KConnected} is the
definition of Section~\ref{sec:prelim}. In \texttt{fullCone B} (whose definition by its adjacency
relation we do not reproduce) the apex \texttt{none} is adjacent to every \texttt{some a}, and
\texttt{some a}, \texttt{some b} are adjacent if and only if $a,b$ are adjacent in $B$.
\texttt{RootedDecomp B r} is $\mathrm{RD}(B,r)$. Theorem~\ref{thm:lovasz-strong} is stated for
\texttt{evenSet G}, the set of vertices of positive even degree.

\subsection{Building and checking}
The artifact of this paper is the directory \texttt{papers/P1-spine/} of the project repository. From
the root of the repository:
{\footnotesize
\begin{verbatim}
export PATH=$HOME/.elan/bin:$PATH
cd lean && lake exe cache get && cd ..   # download the Mathlib build cache
scripts/check_papers.sh P1               # lake build, then P1.lean against axioms.expected
papers/P1-spine/reproduce.sh             # the same, plus the name map and the PDF
\end{verbatim}
}
The entry file \texttt{papers/P1-spine/lean/P1.lean} imports the modules the paper uses. It restates
every headline result with its full statement (Theorem~\ref{thm:main}, both of its directions,
Lemma~\ref{lem:fan}, Theorems~\ref{thm:lovasz-strong}, \ref{thm:lovasz}, \ref{thm:59}
and~\ref{thm:sep}, Propositions~\ref{prop:cone-conn}--\ref{prop:cone-t1}, \ref{prop:cone-leaf}, the
first item of Proposition~\ref{prop:cone-vac}, and Theorem~\ref{thm:cones}(c),(d)) and proves each
by the library declaration. It checks that every other Lean name cited in the paper resolves, and it
runs \texttt{\#print axioms} on each headline result. Its complete output, recorded in
\texttt{axioms.expected} from a clean build on 28 November 2026, consists of sixteen lines of the form
{\footnotesize
\begin{verbatim}
'P1.main_theorem' depends on axioms: [propext, Classical.choice, Quot.sound]
\end{verbatim}
}
one for each restated result; for \texttt{P1.cone\_matchingCut\_iff} the list is \texttt{[propext,
Quot.sound]}. These are the standard axioms of Lean's classical mathematics; no further axiom and no
\texttt{sorryAx} occurs. A full build of the library takes about five minutes on four cores once the
Mathlib cache is downloaded; checking \texttt{P1.lean} against a finished build takes a few seconds.

\subsection{Correspondence}
The table lists each result of the paper with the Lean declarations that prove it (all in namespace
\texttt{Gallai}, files in \texttt{lean/GallaiLean/}). The script
\texttt{papers/P1-spine/scripts/locate.sh} prints the current file and line of every declaration.

{\small
\begin{longtable}{@{}p{0.30\textwidth}p{0.45\textwidth}p{0.20\textwidth}@{}}
\toprule
Paper & Lean declaration & File\\
\midrule
\endhead
Theorem~\ref{thm:main} & \texttt{gallai\_iff\_spineTests'}, \texttt{gallai\_iff\_spineTests} &
Separator\\
Section~\ref{sec:prelim} constructions & \texttt{addLeaf}, \texttt{apex}, \texttt{cone},
\texttt{delVert}, \texttt{delVert\_degree} & Doubling, Realisable, Marked, ThreeConnected\\
Lemma~\ref{lem:kconn} & \texttt{KConnected.le\_degree}, \texttt{KConnected.connected},
\texttt{KConnected.delVert\_connected} & ThreeConnected\\
Lemma~\ref{lem:parity} & \texttt{PathDecomp.ends\_mod\_two}, \texttt{PathDecomp.sum\_ends},
\texttt{PathDecomp.ends\_le\_degree} & Parity, TwoEven\\
Lemma~\ref{lem:split} & \texttt{PathDecomp.exists\_split},
\texttt{PathDecomp.exists\_internal\_of\_ends\_lt}, \texttt{GPath.exists\_join} & Surgery,
ThreeConnected, MarkedSurgery\\
Lemma~\ref{lem:ends} & \texttt{exists\_two\_ends}, \texttt{exists\_three\_ends} & ThreeConnected\\
Lemma~\ref{lem:leaf} & \texttt{leafCost}, \texttt{PathDecomp.trim\_ends}, \texttt{pathNumber\_addLeaf\_le},
\texttt{exists\_leafCost\_le} & PrivateLeaf, ThreeConnected\\
Definition~\ref{def:chain} & \texttt{FanPred}, \texttt{FanAvoid}, \texttt{FanChain} & Fan\\
Lemma~\ref{lem:chain-exec} & \texttt{fan\_exec}, \texttt{PathDecomp.append\_move},
\texttt{PathDecomp.reroute\_move}, \texttt{fanChain\_transport} & Fan, FanMoves\\
Lemma~\ref{lem:chain-exists} & \texttt{fan\_exists\_chain} & Fan\\
Lemma~\ref{lem:fan} & \texttt{fan\_single} & Fan\\
Theorem~\ref{thm:lovasz-strong} & \texttt{lovasz\_ends} & Lovasz\\
Theorem~\ref{thm:lovasz} & \texttt{lovaszStatement} & Lovasz\\
Corollary~\ref{cor:lovasz} & \texttt{lovasz\_tight}, \texttt{lovasz\_on} & SmallBounds\\
Definition~\ref{def:marked} & \texttt{MarkedPartition}, \texttt{markedNumber} & Marked\\
Lemma~\ref{lem:marked-basic} & \texttt{pathNumber\_le\_markedNumber\_add},
\texttt{markedNumber\_anti}, \texttt{markedNumber\_le\_insert},
\texttt{markedNumber\_singleton\_add\_one}, \texttt{exists\_leafCost\_le\_markedNumber} &
Marked, Thm59ClassC, SmallBounds\\
Lemma~\ref{lem:cone} & \texttt{le\_outEnds}, \texttt{apex\_outEnds\_bound} & Marked,
Thm59ClassC\\
Lemma~\ref{lem:hall} & \texttt{markedNumber\_add\_card\_le} & Thm59Hall\\
Definition~\ref{def:MB} & \texttt{MarkedTwoConnTestUpTo}, \texttt{ThreeEndsTestUpTo},
\texttt{MarkedBoundUpTo}, \texttt{SmallBlocks} & MarkedBlocks, ThreeConnected\\
Theorem~\ref{thm:transfer} & \texttt{mbOn\_of\_testsUpTo}, \texttt{markedBoundUpTo\_of\_tests} &
MarkedBlocks, Thm59Transfer\\
Lemma~\ref{lem:glue} & \texttt{glue0}, \texttt{glue1}, \texttt{glueJ}, \texttt{glueU},
\texttt{glueEnd}, \texttt{glueSplit} & MarkedBlocks\\
Lemma~\ref{lem:pendant-edge} & \texttt{PMP.exists\_trim}, \texttt{attach\_pendant} &
MarkedSurgery, MarkedBlocks\\
Steps 0--1 & \texttt{baseStep}, \texttt{twoConnStep} & MarkedBlocks\\
Lemma~\ref{lem:cut} & \texttt{cutStep} & MarkedBlocks\\
Lemma~\ref{lem:leafstep} & \texttt{leafStep} & MarkedBlocks\\
Lemma~\ref{lem:doubleleaf} & \texttt{doubleLeafStep} & MarkedBlocks\\
Lemma~\ref{lem:pendant} & \texttt{pendantStep} & MarkedBlocks\\
Lemma~\ref{lem:residue} & \texttt{residue\_sep} & MarkedBlocks\\
Lemma~\ref{lem:shores} & \texttt{exists\_connected\_matchingCut} & Thm59MatchingCut\\
Lemma~\ref{lem:trimleaves} & \texttt{leafGraph}, \texttt{leafGraph\_trim} & Thm59Leaves\\
Lemma~\ref{lem:lemma6} & \texttt{lemma6} & Thm59Lemma6\\
Theorem~\ref{thm:mc} & \texttt{matchingCut\_markedBound} & Thm59MatchingCut\\
Definition~\ref{def:rooted} & \texttt{RootedAt}, \texttt{RootedTestsUpTo} & ThreeConnected,
Thm59ClassC\\
Theorem~\ref{thm:59} & \texttt{thm59 : Thm59Statement} & Thm59, ThreeConnected\\
Lemma~\ref{lem:coneC} & \texttt{apex\_inClassC} & Thm59ClassC\\
Proposition~\ref{prop:classC} & \texttt{cone\_step}, \texttt{singleton\_step},
\texttt{markedBound\_of\_classC}, \texttt{threeEnds\_of\_classC} & Thm59ClassC\\
Definition~\ref{def:sepcomp} & \texttt{IsSepComponent} & ThreeConnected\\
Theorem~\ref{thm:sep} & \texttt{separatorStatement}, \texttt{separator\_theorem},
\texttt{SeparatorStatement} & Separator, ThreeConnected\\
Section~\ref{sec:sep} setup & \texttt{sideA}, \texttt{sideE} & SeparatorSetup\\
Lemma~\ref{lem:sides} & \texttt{sideE\_induce\_connected}, \texttt{sideE\_port\_pos},
\texttt{sideA\_port\_pos}, \texttt{sideA\_parity}, \texttt{two\_le\_card\_of\_odd} &
SeparatorSetup, SeparatorOdd\\
Lemma~\ref{lem:small}(a)--(c) & \texttt{R2\_on}, \texttt{R3\_on}, \texttt{floor\_on} & SmallBounds\\
Lemma~\ref{lem:root} & \texttt{rootedAt\_of\_lift} & SeparatorSetup\\
Path surgery conventions & \texttt{TPath}, \texttt{exists\_decomp\_of\_edgesM} & PathSurgery\\
Lemma~\ref{lem:surgery} & \texttt{exists\_subst\_list}, \texttt{exists\_cut\_edge},
\texttt{exists\_graft}, \texttt{exists\_graft\_x} & PathSurgery, SeparatorEvenEven\\
Proposition~\ref{prop:odd} & \texttt{rootedAt\_oddComp} & SeparatorOdd\\
Proposition~\ref{prop:evenodd} & \texttt{rootedAt\_evenComp\_oddPorts} & SeparatorEvenOdd\\
Lemma~\ref{lem:local} & \texttt{exists\_localX}, \texttt{exists\_localA'} & SeparatorEvenEven\\
Lemma~\ref{lem:lift-ee} & \texttt{case1\_lift}, \texttt{case2\_lift} (the family $M$, ends at the root);
assembly in \texttt{rootedAt\_evenComp\_evenPorts} & SeparatorEvenEven\\
Proposition~\ref{prop:eveneven} & \texttt{rootedAt\_evenComp\_evenPorts} & SeparatorEvenEven\\
Lemma~\ref{lem:bridge} & \texttt{dbl}, \texttt{two\_mul\_pathNumber\_addLeaf\_le},
\texttt{rootedTest\_of\_gallai} & Doubling\\
Lemma~\ref{lem:gallai-rooted}, Prop.~\ref{prop:forward} & \texttt{test1\_of\_gallai},
\texttt{test2\_of\_gallai}, \texttt{spineTests\_of\_gallai} & ThreeConnected\\
Lemma~\ref{lem:apex} & \texttt{apex\_kConnected\_three} & ThreeConnected\\
Lemma~\ref{lem:fewEven} & \texttt{rootedAt\_of\_few\_even} & ThreeConnected\\
Proposition~\ref{prop:rootedFromTests} & \texttt{rootedAt\_of\_tests},
\texttt{exists\_sepComponent\_or}, \texttt{not\_hasMatchingCut\_delVert\_apex} & ThreeConnected\\
Theorem~\ref{thm:converse} & \texttt{gallai\_of\_spineTestsNMC} & ThreeConnected\\
Proposition~\ref{prop:cone-par} & \texttt{fullCone\_degree\_none}, \texttt{fullCone\_degree\_some},
\texttt{fullCone\_onlyEven\_iff} & EulerianCone\\
Proposition~\ref{prop:cone-conn} & \texttt{kConnected\_three\_fullCone\_iff} & EulerianCone\\
Proposition~\ref{prop:cone-mc} & \texttt{hasMatchingCut\_delVert\_fullCone\_iff},
\texttt{HasMatchingCut.of\_iso} & EulerianCone\\
Proposition~\ref{prop:cone-t1} & \texttt{fullConeDelIso}, \texttt{PathDecomp.exists\_of\_iso},
\texttt{test1At\_fullCone\_iff}, \texttt{test1\_fullCone\_iff} & EulerianCone\\
Proposition~\ref{prop:cone-vac} & \texttt{test2\_fullCone\_apex}, \texttt{test1\_fullCone\_some},
\texttt{test2\_fullCone\_some\_of\_even}, \texttt{not\_unique\_odd} & EulerianCone\\
Proposition~\ref{prop:cone-leaf} & \texttt{rootedDecomp\_iff\_addLeaf} & EulerianCone\\
Theorem~\ref{thm:cones} & \texttt{rootedDecomp\_of\_spineTestsNMC}, \texttt{fullConeTests\_iff},
\texttt{rootedDecomp\_of\_gallai}, \texttt{gallai\_iff\_rootedTest} & EulerianCone, Doubling\\
\bottomrule
\end{longtable}
}

\subsection{How the formal proof is organised}
The formal proof makes some representation choices that the paper does not spell out; none affects
the mathematics.
\begin{itemize}
\item The block transfer (Theorem~\ref{thm:transfer}) runs inside one vertex type: an instance is a
vertex set $U$ carrying all edges of a graph, and sub-instances are restrictions to subsets of $U$,
which are exactly the induced subgraphs used above (\texttt{restrict}, \texttt{SmallerMB},
\texttt{SmallOn}). Its notion of $2$-connectivity (\texttt{TwoConnOn}: at least three vertices,
connected, no cut vertex) is converted to \texttt{KConnected 2} in \texttt{Thm59Transfer.lean}.
\item Private leaves, cones and doubles are built on \texttt{Option V}, \texttt{V ⊕ V} or
\texttt{Option (Bool × V)}, and decompositions are transported along injective vertex maps with their
end counts (\texttt{PathDecomp.project\_full}, \texttt{PathDecomp.map\_injective\_ends},
\texttt{PathDecomp.mapInj}).
\item In the separator theorem all graphs live on the vertex type of $B$, and the exterior graph of
the odd case reuses the vertex $c\in C$ as the leaf; paths are handled as paths of the complete graph
on $V$ with edge multisets (\texttt{TPath}, \texttt{edgesM}, \texttt{endsL}); Section~\ref{sec:sep}
adopts the same convention.
\item In Fan's lemma, a decomposition of $K+tu$ is produced for every graph whose edge set is
$E(K)\cup\{tu\}$, so that the graphs $K_1$ of the induction need not be built as named
constructions.
\end{itemize}

\subsection*{Sources and acknowledgements}
The mathematics of the spine (Sections~\ref{sec:prelim} and~\ref{sec:marked}--\ref{sec:thm114}) is
that of the marked-reduction manuscript~\cite{MR}; the account here was assembled from that
manuscript, from the project's audits of it (A1, A2, A8 and A14), from the Lean formalisation, and
from the earlier write-up~\cite{SP} and its audit A15. Section~\ref{sec:lovasz} is based on the
note~\cite{LF}, and Section~\ref{sec:cones} on Section~2 of~\cite{EC} and its audit A30. The
observation that (T1) on cones is rooted Gallai for the base was first recorded in the project's
synthesis notes. Fan's lemma and its role as an end-transfer move are due to Fan~\cite{Fan05}.

This work was prepared with the assistance of AI tools. All results were machine-checked in Lean~4
with Mathlib, as described in Section~\ref{sec:lean}; the author is responsible for the content.

\begin{thebibliography}{99}
\bibitem{BP19} M.~Bonamy and T.~J.~Perrett, Gallai's path decomposition conjecture for graphs of small
maximum degree, Discrete Math.\ 342 (2019) 1293--1299.
\bibitem{BSCL20} F.~Botler, M.~Sambinelli, R.~S.~Coelho and O.~Lee, Gallai's path decomposition
conjecture for graphs with treewidth at most 3, J.\ Graph Theory 93 (2020) 328--349.
\bibitem{BS21} F.~Botler and M.~Sambinelli, Towards Gallai's path decomposition conjecture, J.\ Graph
Theory 97 (2021) 161--184.
\bibitem{Fan05} G.~Fan, Path decompositions and Gallai's conjecture, J.\ Combin.\ Theory Ser.~B 93
(2005) 117--125.
\bibitem{Hall35} P.~Hall, On representatives of subsets, J.\ London Math.\ Soc.\ 10 (1935) 26--30.
\bibitem{L68} L.~Lov\'asz, On covering of graphs, in: Theory of Graphs (Proc.\ Colloq.\ Tihany,
1966), Academic Press, New York, 1968, 231--236.
\bibitem{Mathlib} The mathlib Community, The Lean mathematical library, in: Proc.\ 9th ACM SIGPLAN
Int.\ Conf.\ on Certified Programs and Proofs (CPP 2020), 367--381.
\bibitem{Lean4} L.~de Moura and S.~Ullrich, The Lean 4 theorem prover and programming language, in:
Automated Deduction -- CADE 28, Lecture Notes in Comput.\ Sci.\ 12699, Springer, 2021, 625--635.
\bibitem{Pyb96} L.~Pyber, Covering the edges of a connected graph by paths, J.\ Combin.\ Theory
Ser.~B 66 (1996) 152--159.
\bibitem{MR} [Author], Cut vertices and cycle intersections in Gallai's path decomposition problem,
unrefereed manuscript, 24 November 2026; \texttt{gallai/01-original-manuscripts/marked-reduction.pdf}
in the project repository. Theorems 1, 7, 59, 113 and 114, Lemma~6, \S1, \S2 and \S17.
\bibitem{SP} [Author], The reduction spine for Gallai's path decomposition conjecture: two deletion
tests on $3$-connected graphs, working note, November 2026; \texttt{manuscripts/spine} in the project
repository. Superseded by the present paper.
\bibitem{LF} [Author], Lov\'asz's path decomposition theorem from Fan's addible-edge lemma, working
note, November 2026; \texttt{manuscripts/lovasz-fan} in the project repository. Superseded by
Section~\ref{sec:lovasz}.
\bibitem{EC} [Author], The two-end test on cones: rooted Gallai for Eulerian graphs, working note,
November 2026; \texttt{manuscripts/eulerian-cone} in the project repository. Section~2 is superseded
by Section~\ref{sec:cones}.
\bibitem{P2} [Author], Constrained Lov\'asz decompositions: separation, neighbour avoidance, leaf
tails and prescribed ends, in preparation (paper P2 of the project).
\bibitem{P3} [Author], Reversible lifts from constrained Lov\'asz decompositions to the spine tests,
in preparation (paper P3 of the project).
\bibitem{P4} [Author], Rooted Gallai for Eulerian graphs, and the two-end test at near-universal
roots, in preparation (paper P4 of the project).
\end{thebibliography}

\end{document}
